\documentclass[11pt,a4paper]{article}
\usepackage[T1]{fontenc}
\usepackage[utf8]{inputenc}
\usepackage{lmodern}
% La vía moderna de babel permite importar los rótulos españoles también
% en instalaciones que no incluyen el archivo tradicional spanish.ldf.
\IfFileExists{spanish.ldf}{%
  \usepackage[spanish,es-nodecimaldot,es-noquoting]{babel}%
}{%
  \usepackage{babel}
  \babelprovide[import,main]{spanish}%
}
\usepackage{amsmath,amssymb,amsthm}
\usepackage{mathtools}
\usepackage[margin=25mm,headheight=15pt,headsep=10mm,footskip=12mm]{geometry}
\usepackage{microtype}
\usepackage{xcolor}
\usepackage{booktabs,tabularx}
\usepackage{enumitem}
\usepackage{fancyhdr}
\usepackage{hyperref}
\usepackage{bookmark}

\definecolor{azulnota}{HTML}{203B55}
\definecolor{grisnota}{HTML}{59636D}
\hypersetup{
  colorlinks=true,
  linkcolor=azulnota,
  citecolor=azulnota,
  urlcolor=azulnota,
  pdftitle={Hipótesis de Riemann — Notas de estudio},
  pdfauthor={Alevi Peña},
  pdfsubject={Lectura de dos exposiciones del argumento de Alpöge para las funciones L de Dirichlet},
  pdfkeywords={Riemann, Dirichlet, Mellin, retículos, notas de estudio}
}
\pagestyle{fancy}
\fancyhf{}
\fancyhead[L]{\small\textcolor{grisnota}{Hipótesis de Riemann}}
\fancyhead[R]{\small\textcolor{grisnota}{Notas de estudio}}
\fancyfoot[C]{\small\thepage}
\renewcommand{\headrulewidth}{0.3pt}
\setlength{\parindent}{0pt}
\setlength{\parskip}{0.55em}
\setlength{\emergencystretch}{2em}
\setlist{topsep=0.4em,itemsep=0.3em,parsep=0pt}
\setlist[itemize,enumerate]{leftmargin=*}
\setlist[description]{style=nextline,leftmargin=1.4em,labelindent=0pt}
\newcolumntype{Y}{>{\raggedright\arraybackslash}X}
\setcounter{tocdepth}{2}
\numberwithin{equation}{section}
\allowdisplaybreaks[1]
\widowpenalty=10000
\clubpenalty=10000

\theoremstyle{plain}
\newtheorem{proposicion}{Proposición}[section]
\newtheorem{teorema}[proposicion]{Teorema}
\newtheorem{lema}[proposicion]{Lema}
\theoremstyle{definition}
\newtheorem{definicion}[proposicion]{Definición}
\newtheorem{ejercicio}{Ejercicio}[section]
\theoremstyle{remark}
\newtheorem*{observacion}{Observación}
\renewcommand{\proofname}{Demostración}
\renewcommand{\contentsname}{Índice}
\renewcommand{\refname}{Fuentes de esta entrega}

\newcommand{\N}{\mathbb{N}}
\newcommand{\Z}{\mathbb{Z}}
\newcommand{\Q}{\mathbb{Q}}
\newcommand{\R}{\mathbb{R}}
\newcommand{\C}{\mathbb{C}}
\newcommand{\OK}{\mathcal{O}_K}
\newcommand{\dd}{\,\mathrm{d}}
\DeclareMathOperator{\Ree}{Re}
\DeclareMathOperator{\Res}{Res}
\DeclareMathOperator{\Hom}{Hom}
\DeclareMathOperator{\length}{length}
\DeclareMathOperator{\ord}{ord}
\newcommand{\fuente}[1]{\par\small\textcolor{grisnota}{\textbf{Lectura de referencia.} #1}\par\normalsize}


\title{\textbf{Hipótesis de Riemann}\\[0.35em]
  \Large Notas de estudio}
\author{Alevi Peña}
\date{22 de noviembre de 2026}

\begin{document}
\maketitle
\thispagestyle{empty}
\begin{center}
  \small Cuaderno de lectura y desarrollo matemático
\end{center}
\vspace{0.5em}
\noindent Estas notas organizan el estudio de dos exposiciones de la propuesta
de Alpöge para las funciones $L$ de Dirichlet. La introducción presenta
los primos, zeta, los caracteres, los enunciados de RH y GRH y el vocabulario
analítico y aritmético utilizado en el cuaderno. Las notas desarrollan la relación entre
sumas sobre potencias de primos y polos de una transformada, y demuestran
cómo una cota superior en ventanas permite excluir esos polos.
Incluyen también un estudio de las proyecciones racionales y los
denominadores de retículos. La identidad de caracteres enlaza las
condiciones de Frobenius con las sumas corregidas y explica la selección
de los coeficientes positivos. La Nota~\ref{sec:marcos} demuestra la
desigualdad de marcos y su integración para transferir pendientes medias.

\tableofcontents
\clearpage

\setcounter{section}{-1}
\section{Guía del cuaderno}
\label{sec:guia}

\subsection{Fuentes y modo de lectura}

Llamaremos \textbf{A} a \emph{Dirichlet L-functions and the critical line},
el documento de 55 páginas, y \textbf{B} a la nota preliminar de Robin Zhang,
de 73 páginas. Las páginas citadas corresponden a los ejemplares suministrados;
sus números impresos coinciden con las páginas del PDF. Las referencias
completas se encuentran al final del cuaderno \cite{fuenteA,fuenteB}.

B presenta los objetos de manera más progresiva y explica antecedentes
matemáticos. A resulta especialmente útil para localizar hipótesis,
normalizaciones, estimaciones uniformes y nombres de declaraciones Lean.
Los utilizaremos conjuntamente: primero la motivación y después el detalle
necesario para reconstruir un paso.

El alcance que estudiamos es la hipótesis de Riemann generalizada para
\emph{funciones $L$ de Dirichlet}. La función zeta corresponde al carácter
primitivo de conductor $1$. Las formas automorfas intervienen como objetos
auxiliares; su presencia no convierte el enunciado en una afirmación sobre
todas las funciones $L$ automorfas. Véanse A, p.~1, y B, p.~3.

En estas notas distinguimos tres estados de trabajo:
\begin{description}
  \item[Desarrollado aquí.] Se incluye el cálculo o la demostración del paso.
  \item[Resultado de la fuente.] Se identifica el enunciado que utilizamos,
    sin sustituir su prueba por una explicación informal.
  \item[Pendiente de estudio.] Falta reconstruir el argumento o cotejarlo
    con mayor detalle. Esto describe el estado del cuaderno.
\end{description}

La nota B advierte expresamente, en p.~2, que su elaboración no permite
afirmar una traducción completamente fiel de todo el código. Nuestros
cotejos entre A y B tampoco constituyen una auditoría de ese código.

\subsection{La arquitectura en seis bloques}

\begin{enumerate}[label=\textbf{\arabic*.}]
  \item \textbf{Detectar ceros con una transformada.}
    La derivada logarítmica de $L$ transforma cada cero en un polo.
    Una suma suavizada sobre potencias de primos tiene una transformada de
    Mellin que conserva esos polos. La corrección en $s=1$ elimina un término
    ajeno a los ceros buscados. Esta es la Nota~\ref{sec:sumas}.
  \item \textbf{Extraer los caracteres adecuados.}
    La ortogonalidad de caracteres de un grupo de Galois abeliano expresa
    condiciones sobre Frobenius como sumas de caracteres. Los pesos y
    promedios seleccionados conservan una combinación positiva de las
    funciones corregidas. Referencias: A, §4; B, §3.
  \item \textbf{Comparar divisibilidad y normas de periodos.}
    Los periodos automorfos normalizados dan aplicaciones entre retículos.
    La divisibilidad permite ampliar los dominios sin perder la integridad
    de sus imágenes; las normas complejas controlan la altura de las
    aplicaciones. Referencias: A, §6; B, §4.1 y §4.5.
  \item \textbf{Transferir pendientes medias.}
    Una familia de mapas debe cubrir suficientes direcciones del espacio
    destino. Las estimaciones de marcos, junto con la filtración de
    Harder--Narasimhan, convierten ese control en una desigualdad para
    pendientes medias. Referencias: A, §5; B, §4.4.
  \pagebreak
  \item \textbf{Controlar los denominadores.}
    La descomposición en autoespacios puede romper la estructura integral.
    Las congruencias de Hecke y las extensiones de representaciones de Galois
    permiten medir el coste mediante grupos cohomológicos. La
    Nota~\ref{sec:proyecciones} desarrolla el ejemplo elemental de entrada.
  \item \textbf{Elegir escalas y una ventana común.}
    Después de cancelar los términos principales, aparecen costes del orden
    de $x/m$ y de $m$, salvo potencias pequeñas de $x$. Una longitud de torre
    $m$ cercana a $\sqrt{x}$ los equilibra. Las condiciones deben cumplirse
    en un mismo intervalo. Referencias: A, §§10--11; B, §§5.5--6.
\end{enumerate}

\subsection{Notación común y primeras metas}

Reservamos $G_\chi$ para la suma corregida de un carácter y $H_\psi$ para la
combinación positiva de las potencias de $\psi$. En B esta última se llama
$\mathcal G_\psi$. Usaremos $o$ para el orden del carácter, a fin de no
confundirlo con rangos de espacios automorfos. Las letras $m$ y $p$ designan,
respectivamente, una longitud de torre o una multiplicidad según se indique,
y un primo.

\begin{center}
\small
\begin{tabularx}{\textwidth}{@{}>{\raggedright\arraybackslash}p{0.34\textwidth}YY@{}}
\toprule
\textbf{Paso que queremos reconstruir} & \textbf{A} & \textbf{B} \\
\midrule
Transformada y corrección & §2, pp.~3--7 & §2, pp.~14--19 \\
Identidad de caracteres & Prop.~4.2, p.~15 & Prop.~3.1, p.~21 \\
Marcos y pendientes & Prop.~5.2 y Teo.~5.3, pp.~18--20 & Lema~4.2 y Prop.~4.3, pp.~41--44 \\
Proyecciones y denominadores & §§6--8 & Lema~5.1, p.~53; pp.~54--56 \\
Cancelación y raíz cuadrada & Props.~10.1--10.2, pp.~47--49 & (4.5)--(4.10), pp.~48--51; Lema~6.1, pp.~65--66 \\
\bottomrule
\end{tabularx}
\end{center}

La meta inmediata es poder demostrar los cálculos de las dos notas siguientes
sin consultar el texto. El siguiente desarrollo será el criterio unilateral
de ventanas: justificar la concentración del núcleo y explicar cómo se
excluye un polo a partir de la transformada. Después estudiaremos la identidad
de caracteres y la desigualdad de marcos. Las construcciones geométricas y cohomológicas
permanecen identificadas como resultados de las fuentes hasta desarrollar
sus hipótesis y demostraciones.

\clearpage
\section{Sumas sobre primos y transformadas}
\label{sec:sumas}

\fuente{A, §2, especialmente Prop.~2.1, pp.~3--4; B, §1.2,
pp.~3--5, y §2, pp.~14--19. Los cálculos de esta nota se desarrollan aquí;
el criterio completo de ventanas se demuestra en la Nota~\ref{sec:ventanas}.}

\subsection{La suma suavizada}

Sea $\chi$ un carácter primitivo de Dirichlet. La función de von Mangoldt es
\[
\Lambda(n)=
\begin{cases}
\log p,& n=p^a,\quad p\text{ primo},\quad a\ge1,\\
0,&\text{en otro caso}.
\end{cases}
\]
Para $x\ge1$ definimos
\begin{equation}
  S_\chi(x)=\sum_{n\le x}\Lambda(n)\chi(n)\left(\frac{x}{n}-1\right).
  \label{eq:suma-suave}
\end{equation}
La suma es finita. En $x=n$, el término que acaba de incorporarse vale cero:
por eso el peso suaviza los saltos de una suma de conteo. El soporte de
$\Lambda$ hace que solo contribuyan potencias de primos.

En el semiplano $\Ree s>1$, la derivación logarítmica del producto de Euler
proporciona la identidad clásica
\begin{equation}
  -\frac{L'(s,\chi)}{L(s,\chi)}
   =\sum_{n\ge1}\frac{\Lambda(n)\chi(n)}{n^s}.
  \label{eq:derivada-log}
\end{equation}
Allí las series convergen absolutamente. No se está suponiendo información
sobre ceros dentro de la franja crítica.

\subsection{Cálculo de la transformada de Mellin}

\begin{proposicion}
Para $\Ree s>1$,
\begin{equation}
  \int_1^\infty S_\chi(x)x^{-s-1}\dd x
   =-\frac{L'(s,\chi)}{s(s-1)L(s,\chi)}.
  \label{eq:mellin-s}
\end{equation}
\end{proposicion}

\begin{proof}
Escribamos $\sigma=\Ree s>1$. Para un entero $n\ge1$,
\begin{align*}
\int_n^\infty\left(\frac{x}{n}-1\right)x^{-s-1}\dd x
 &=\frac1n\int_n^\infty x^{-s}\dd x
      -\int_n^\infty x^{-s-1}\dd x\\
 &=\frac{n^{-s}}{s-1}-\frac{n^{-s}}s
  =\frac{n^{-s}}{s(s-1)}.
\end{align*}
Antes de intercambiar suma e integral comprobamos convergencia absoluta.
Como $|\chi(n)|\le1$ y $\Lambda(n)\le\log n$,
\begin{align*}
&\sum_{n\ge1}\Lambda(n)|\chi(n)|
  \int_n^\infty\left(\frac{x}{n}-1\right)x^{-\sigma-1}\dd x\\
&\hspace{2em}\le
 \frac{1}{\sigma(\sigma-1)}
 \sum_{n\ge1}\frac{\log n}{n^\sigma}<\infty.
\end{align*}
Fubini permite el intercambio. El cálculo del núcleo y
\eqref{eq:derivada-log} dan \eqref{eq:mellin-s}.
\end{proof}

La razón para elegir $x/n-1$ ya es visible: su transformada introduce el
factor explícito $1/[s(s-1)]$. Un cero de $L$ se convierte en un polo de
$L'/L$, y podremos seguirlo a través de este factor.

\subsection{La corrección en el punto \texorpdfstring{$s=1$}{s=1}}

Supongamos primero que $\chi$ no es principal. Utilizamos el resultado
clásico $L(1,\chi)\ne0$ y ponemos
\[
  \ell_\chi=\frac{L'(1,\chi)}{L(1,\chi)},\qquad
  G_\chi(x)=S_\chi(x)+x\ell_\chi.
\]
Para $\Ree s>1$, la transformada corregida es
\begin{align}
F_\chi(s)
  &:=\int_1^\infty G_\chi(x)x^{-s-1}\dd x\notag\\
  &=-\frac{L'(s,\chi)}{s(s-1)L(s,\chi)}
      +\frac{\ell_\chi}{s-1}
   =\frac{-L'(s,\chi)/L(s,\chi)+s\ell_\chi}{s(s-1)}.
   \label{eq:mellin-g}
\end{align}
El numerador se anula en $s=1$. Como $L'/L$ es holomorfa cerca de ese
punto, la singularidad es removible. \textbf{El signo positivo de
$x\ell_\chi$ es indispensable}: con el signo contrario no se produciría
esta cancelación.

Para el carácter primitivo principal, $L(s,1)=\zeta(s)$ y la corrección cambia:
\begin{equation}
 G_1(x)=S_1(x)-x\log x+(1+\gamma)x,
 \label{eq:g-principal}
\end{equation}
donde $\gamma$ es la constante de Euler. Por integración directa,
\begin{equation}
 F_1(s)=
 -\frac{\zeta'(s)}{s(s-1)\zeta(s)}
 -\frac1{(s-1)^2}+\frac{1+\gamma}{s-1}.
 \label{eq:mellin-principal}
\end{equation}
Para comprobar la cancelación, escribimos $t=s-1$ y usamos
\[
 \zeta(1+t)=\frac1t+\gamma+O(t),\qquad
 -\frac{\zeta'(1+t)}{\zeta(1+t)}=\frac1t-\gamma+O(t).
\]
El primer término de \eqref{eq:mellin-principal} es
$t^{-2}-(1+\gamma)t^{-1}+O(1)$. Los otros dos lo cancelan exactamente
hasta dejar una función holomorfa en $s=1$.

\subsection{El residuo en un cero}

\begin{proposicion}
Sea $\rho$ un cero de multiplicidad $q\ge1$ de $L(s,\chi)$, con
$0<\Ree\rho<1$. La continuación meromorfa de $F_\chi$ tiene en $\rho$
un polo simple cuyo residuo es
\begin{equation}
 \Res_{s=\rho}F_\chi(s)
 =-\frac{q}{\rho(\rho-1)}
 =\frac{q}{\rho(1-\rho)}.
 \label{eq:residuo}
\end{equation}
\end{proposicion}

\begin{proof}
Localmente, $L(s,\chi)=(s-\rho)^q u(s)$, donde $u$ es holomorfa y
$u(\rho)\ne0$. Por tanto
\[
 \frac{L'(s,\chi)}{L(s,\chi)}
 =\frac{q}{s-\rho}+\frac{u'(s)}{u(s)}.
\]
El factor $s(s-1)$ no se anula en $\rho$, y los términos de corrección
son holomorfos allí. La fórmula del residuo se sigue de
\eqref{eq:mellin-g} o \eqref{eq:mellin-principal}.
\end{proof}

Este argumento no sitúa los ceros en la recta crítica. Describe qué huella
dejaría en la transformada cualquier cero de la franja, incluso uno con
parte real mayor que $1/2$.

\subsection{Positividad y el objetivo de las ventanas}

Si $\psi$ tiene orden $o$, denotamos por $\psi^{[j]}$ el carácter primitivo
que induce $\psi^j$ y definimos
\begin{equation}
 H_\psi(x)=G_1(x)+\frac1{3o+1}
    \sum_{j=1}^{o-1}\Ree G_{\psi^{[j]}}(x).
 \label{eq:hpsi}
\end{equation}
Para $o=1$ la suma es vacía. Los caracteres retenidos forman una familia
estable bajo conjugación, de modo que la suma de sus partes reales coincide
con la suma de las funciones. En un cero compartido, todos los residuos
son múltiplos de $1/[\rho(1-\rho)]$ con multiplicidades y coeficientes
positivos. Su suma no puede cancelarse.

El objetivo aritmético puede formularse así: para cada $\eta>0$ y cada
cota inicial $A\ge2$, existe $Y\ge A$ tal que
\begin{equation}
 H_\psi(x)\le x^{1/2+\eta}
 \qquad\text{para todo }x\in[Y,Y^{3/2}].
 \label{eq:ventana}
\end{equation}
Es una cota unilateral en intervalos arbitrariamente lejanos. No se afirma
una cota de valor absoluto ni que los intervalos cubran una semirrecta.
El criterio analítico está adaptado a esos cuantificadores. La
Nota~\ref{sec:ventanas} desarrolla su demostración mediante derivadas de
la transformada de Laplace, concentración de un núcleo y selección de
un polo en discos compactos. La construcción aritmética de las ventanas
sigue pendiente de estudio. Véanse A, §2, y B, Props.~2.1--2.3.

\begin{ejercicio}
Compruebe directamente que
$\int_1^\infty x^{-s}\dd x=(s-1)^{-1}$ y
$\int_1^\infty(\log x)x^{-s}\dd x=(s-1)^{-2}$ para $\Ree s>1$.
Explique por qué estas integrales no se usan como integrales convergentes
en toda la franja crítica: allí se utiliza la continuación meromorfa de
las expresiones obtenidas.
\end{ejercicio}

\clearpage
\section{Proyecciones y denominadores}
\label{sec:proyecciones}

\fuente{B, §5.1, especialmente Lema~5.1, p.~53, y §§5.2--5.3,
pp.~54--57; A, Lema~6.3, p.~22, y §8, pp.~35--42. El ejemplo de
dimensión dos y el lema local se demuestran aquí.}

\subsection{Un operador integral con un proyector no integral}

Sea $p$ primo, sea $V=\Q^2$ y sea $M=\Z^2$ su retículo habitual.
Consideremos
\begin{equation}
 T=\begin{pmatrix}0&1\\0&p\end{pmatrix},
 \qquad T(x,y)=(y,py).
 \label{eq:operador-t}
\end{equation}
La matriz tiene entradas enteras, así que $T(M)\subseteq M$. Sus dos
autovalores son $0$ y $p$, con autoespacios
\[
 V_0=\Q(1,0),\qquad V_p=\Q(1,p).
\]
Todo vector se descompone de manera única como
\begin{equation}
 (x,y)=\left(x-\frac yp\right)(1,0)+\frac yp(1,p).
 \label{eq:descomposicion}
\end{equation}
Por tanto, el proyector sobre $V_p$ a lo largo de $V_0$ es
\begin{equation}
 P=\frac{T}{p}
   =\begin{pmatrix}0&1/p\\0&1\end{pmatrix}.
 \label{eq:proyector-p}
\end{equation}
La igualdad $T^2=pT$ comprueba que $P^2=P$. Sin embargo,
$P(0,1)=(1/p,1)\notin M$. El operador $T$ es integral y diagonalizable
sobre $\Q$, pero su proyector espectral no preserva $M$.

La diferencia entre los autovalores es $p$, que no es una unidad sobre
$\Z_p$. La fórmula $P=(T-0I)/(p-0)$ muestra exactamente dónde aparece
el denominador. Módulo $p$, los autovalores coinciden y
\[
 \overline T=\begin{pmatrix}0&1\\0&0\end{pmatrix}
\]
es un operador nilpotente no nulo. La separación de autoespacios sobre
el cuerpo no se conserva como descomposición del retículo.

\subsection{Dos cocientes que conviene distinguir}

Los retículos que están dentro de cada autoespacio son
\[
 M\cap V_0=\Z(1,0),\qquad M\cap V_p=\Z(1,p).
\]
Su suma es
\[
 M'=(M\cap V_0)\oplus(M\cap V_p)
    =\{(a+b,pb):a,b\in\Z\}=\Z\times p\Z.
\]
Así, $[M:M']=p$. Este índice mide el coste de sustituir $M$ por un
retículo que sí se descompone en los autoespacios.

El proyector, en cambio, produce
\[
 PM=\frac1p\Z(1,p),
 \qquad
 \frac{PM}{M\cap V_p}
 =\frac{(1/p)\Z(1,p)}{\Z(1,p)}\simeq\Z/p\Z.
\]
Este segundo cociente mide cuánto amplía el proyector el retículo de
$V_p$. Es el objeto que aparece directamente en el Lema~5.1 de B.
Ambos índices valen $p$ en el ejemplo, pero sus definiciones son distintas.

\subsection{El lema local de denominadores}

\begin{lema}
Sea $\mathcal O$ un anillo de valoración discreta con cuerpo de fracciones
$K$ y valoración $v$ normalizada por $v(\varpi)=1$. Sea $M$ un retículo
pleno en $V=V_1\oplus V_2$, es decir, un $\mathcal O$-submódulo
finitamente generado que genera $V$ sobre $K$, y sea $e$ el proyector sobre $V_1$.
Supongamos que un operador $T$ preserva $M$, actúa por un escalar
$a\in\mathcal O\setminus\{0\}$ sobre $V_1$ y por cero sobre $V_2$.
Entonces
\begin{equation}
 aeM\subseteq M\cap V_1,
 \qquad
 \length_{\mathcal O}\frac{eM}{M\cap V_1}
 \le (\dim_K V_1)v(a).
 \label{eq:lema-denominadores}
\end{equation}
\end{lema}

\begin{proof}
Se tiene $T=ae$ en $V$. Para $m\in M$, el vector $aem=Tm$ pertenece
a $M$ y a $V_1$. Además, $M\cap V_1\subseteq eM$, pues $e$ es la
identidad en $V_1$. El cociente de \eqref{eq:lema-denominadores} es,
por tanto, un cociente de $eM/aeM$.

El módulo $eM$ es finitamente generado, libre de torsión y genera $V_1$;
sobre un anillo de valoración discreta es libre de rango $d_1=\dim_K V_1$.
Por ello
\[
 \length_{\mathcal O}(eM/aeM)
 =d_1\length_{\mathcal O}(\mathcal O/a\mathcal O)=d_1v(a).
\]
La longitud no aumenta al pasar a un cociente, lo que prueba la cota.
\end{proof}

Para \eqref{eq:operador-t}, después de extender de $\Z$ a $\Z_p$, se
tiene $a=p$, $d_1=1$ y longitud $1$. La desigualdad es una igualdad.
En una valoración discreta, \emph{longitud} cuenta factores isomorfos al
cuerpo residual; sobre $\Z_p$, un módulo finito de longitud $r$ tiene
orden $p^r$.

\subsection{Del ejemplo a los periodos automorfos}

Los espacios automorfos utilizados por las fuentes admiten descomposiciones
en sistemas de autovalores de Hecke sobre un cuerpo de coeficientes.
Los operadores de Hecke preservan estructuras integrales, pero proyectar
en los autoespacios elegidos puede introducir denominadores. Si dos
sistemas de autovalores son congruentes módulo un primo de coeficientes,
un operador que anula uno puede actuar por un escalar divisible por ese
primo sobre el otro. El lema cuantifica el coste de dividir por tal escalar.

La conexión con la comparación de pendientes puede resumirse sin ocultar
sus objetos. Para un cuerpo de coeficientes $K$, las fuentes construyen
un emparejamiento normalizado
\[
 P_\alpha:U_\alpha^\circ\times W\longrightarrow K
\]
y retículos $M'_\alpha\subset U_\alpha^\circ$, $M_W\subset W$.
La divisibilidad da un ideal fraccionario $\mathfrak a_\alpha$ con
\[
 P_\alpha(M'_\alpha,M_W)\subseteq\mathfrak a_\alpha.
\]
Al definir
\[
 N_\alpha=\mathfrak a_\alpha^{-1}M'_\alpha,
 \qquad C=\Hom_{\OK}(M_W,\OK),
\]
se obtiene una aplicación integral
\[
 b_\alpha:N_\alpha\longrightarrow C,
 \qquad b_\alpha(f)(g)=P_\alpha(f,g).
\]
Ampliar el retículo fuente consume la divisibilidad disponible; estimar
las normas de $b_\alpha$ permite comparar esa ganancia con el coste
complejo. Los denominadores de las proyecciones restan parte de la ganancia
y deben entrar en la estimación. Véanse B, pp.~29 y 49--56, y A,
Lema~6.3 y §8.

\subsection{Por qué el siguiente paso es cohomológico}

En la construcción profunda de las fuentes, las congruencias entre formas
automorfas se estudian mediante extensiones de sus representaciones de
Galois. Las condiciones locales de esas extensiones llevan, entre otros
objetos, a los grupos
\begin{equation}
 S_\ell(E,j)=\ker\left(
 H^1(G_{E,S_\ell},\Q_\ell/\Z_\ell(-j))
 \longrightarrow\prod_{w\mid\ell}
 H^1(E_w,\Q_\ell/\Z_\ell(-j))\right),
 \label{eq:grupo-selmer}
\end{equation}
donde $j\ge1$ y $S_\ell$ es el conjunto de lugares de $E$ sobre $\ell$.
El signo del twist es parte de la definición. Si el grupo es finito, su
longitud multiplicada por $\log\ell$ es el logaritmo de su orden.

Para el argumento no basta saber que un grupo aislado es finito. Se
necesita controlar sumas de estas longitudes sobre primos y twists, con
uniformidad al variar los campos auxiliares. A, §8, y B, §5, formulan
esas cotas y la selección de una ventana común. Esta nota explica por
qué hacen falta; su construcción sigue siendo un resultado de las fuentes
pendiente de desarrollar en el cuaderno.

\begin{ejercicio}
Sustituya $p$ por un entero $d\ge2$ en \eqref{eq:operador-t}.
Calcule $[M:M']$ y $PM/(M\cap V_d)$. Después localice en un primo
$\ell$: ¿qué longitud obtiene y para qué primos el proyector ya es integral?
\emph{Comprobación:} el índice es $d$, el cociente es $\Z/d\Z$ y la
longitud local es $v_\ell(d)$.
\end{ejercicio}

\clearpage
\section{Criterio unilateral de ventanas}
\label{sec:ventanas}

Esta nota completa el paso analítico que quedó abierto en la
Nota~\ref{sec:sumas}. Demostramos por qué una cota superior sobre intervalos
completos, disponibles tan lejos como se quiera, basta para excluir polos
de una transformada. La prueba conserva el signo de la función, controla
las colas de un núcleo y selecciona un polo mediante distancias en discos
compactos.

\fuente{A, Teorema~2.2, Lemas~2.3--2.4 y Corolario~2.5, pp.~4--6;
B, Proposición~2.1, Lema~2.2 y demostración, pp.~14--18. Al final se
documenta el cotejo con declaraciones Lean en una revisión fijada.
La demostración se reconstruye en lenguaje matemático; no se presenta
como una nueva ejecución de Lean.}

\subsection{Enunciado y orden de los cuantificadores}

Sea $g:[0,\infty)\to\R$ medible y supongamos que existe $C_0\ge0$ tal que
\begin{equation}
 |g(u)|\le C_0e^{2u}\qquad(u\ge0).
 \label{vent-crecimiento}
\end{equation}
Su transformada de Laplace, definida inicialmente para $\Ree s>2$, es
\[
 \Phi(s)=\int_0^\infty g(u)e^{-su}\dd u.
\]
Suponemos que admite continuación meromorfa a $\Ree s>0$, que todos sus
polos son simples y que el residuo de cada polo real es un número real
positivo. Denotamos por $Z$ su conjunto de polos. La meromorfía implica
que cada compacto contenido en ese semiplano contiene finitos polos.

Fijemos $0<\theta<1$ y $\delta>0$. La \emph{hipótesis integrada de ventanas}
significa lo siguiente: para cada $\eta>0$ y cada $U_0\ge1$, existe
$U\ge U_0$ y una función $R_U$ sobre $W_U=[U,(1+\delta)U]$ tales que
\[
 R_U(u)\ge0,\qquad g(u)\le R_U(u)\quad(u\in W_U),
\]
el producto $R_U(u)e^{-\theta u}$ es integrable en $W_U$ y
\begin{equation}
 \int_U^{(1+\delta)U}R_U(u)e^{-\theta u}\dd u\le e^{\eta U}.
 \label{vent-hipotesis}
\end{equation}
La ventana y su mayorante pueden depender de $\eta$ y del umbral.
Los parámetros $\theta,\delta$ y la constante $C_0$ permanecen fijos.
Exigir $U_0\ge1$ es equivalente a permitir cualquier umbral real: basta
reemplazarlo por $\max\{1,U_0\}$. Las integrales de esta nota son integrales
de Lebesgue; sus valores no cambian al añadir o quitar extremos de intervalos.

\begin{proposicion}[Criterio unilateral]
\label{vent-criterio}
Bajo estas hipótesis, todo polo $\rho\in Z$ satisface
$\Ree\rho\le\theta$.
\end{proposicion}

Esta formulación sigue la interfaz general \texttt{PoleData} junto con
\texttt{WindowHypU} de Lean. A añade la localización de los polos en
$\Ree s\le1$ y la regularidad en $s=1$, propias de su aplicación.
Nuestra prueba solo necesita que los polos queden a la izquierda de un
centro $\sigma>2$; esto se deduce de la propia representación de Laplace.
B parte directamente de una cota puntual, que obtendremos como corolario.

\subsection{Derivadas normalizadas y la desigualdad unilateral}

Para $a>0$ y todo entero $k\ge0$, integración por partes da
\begin{equation}
 \int_0^\infty\frac{u^ke^{-au}}{k!}\dd u=a^{-k-1}.
 \label{vent-integral-gamma}
\end{equation}
En efecto, la integral sin el factorial vale $1/a$ si $k=0$, y la
recurrencia es $I_k=(k/a)I_{k-1}$. Los términos de borde se anulan
porque la exponencial decreciente domina cualquier potencia.

Si $\Ree z_0>2$, ponemos $b_0=(\Ree z_0-2)/2>0$.
En el disco $|z-z_0|<b_0$, cada derivada del integrando de orden $j$
tiene módulo a lo sumo $C_0u^je^{-b_0u}$, integrable por
\eqref{vent-integral-gamma}. Esta dominación, también para el orden
siguiente, justifica diferenciar bajo el signo integral. Por tanto,
$\Phi$ es holomorfa en $\Ree s>2$ y
\begin{equation}
 A_k(z):=\frac{(-1)^k}{k!}\Phi^{(k)}(z)
 =\int_0^\infty g(u)\frac{u^ke^{-zu}}{k!}\dd u.
 \label{vent-derivadas}
\end{equation}
En particular, no hay polos con parte real mayor que $2$.
Todo polo queda estrictamente a la izquierda de cualquier $\sigma>2$.
Las identidades integrales se utilizan en este semiplano; más a la
izquierda se utiliza la continuación meromorfa.

Como $g$ es real, $A_k(\sigma)$ es real. Escribiendo
$g_+(u)=\max\{g(u),0\}$, la desigualdad triangular da
\begin{equation}
 \begin{aligned}
 |A_k(\sigma+it)|+A_k(\sigma)
 &\le\int_0^\infty (|g(u)|+g(u))
           \frac{u^ke^{-\sigma u}}{k!}\dd u\\
 &=2\int_0^\infty g_+(u)\frac{u^ke^{-\sigma u}}{k!}\dd u.
 \end{aligned}
 \label{vent-unilateral}
\end{equation}
La segunda contribución del lado izquierdo conserva su signo. Esta
identidad permite utilizar un mayorante de $g$ en la ventana sin pedir
allí una mejora de la cota para $|g|$. Fuera de la ventana sigue disponible
el crecimiento elemental \eqref{vent-crecimiento}.

{\raggedright
Véanse A, (2.9)--(2.10), p.~5; B, p.~17; y las declaraciones
\texttt{iteratedDeriv\_laplace} y \texttt{landau\_window\_ineq}.
\par}

\subsection{Concentrar el núcleo en una ventana completa}

\begin{lema}[Control uniforme de las colas]
\label{vent-colas}
Para cada $\delta>0$ existen $\sigma>2$ y $C_\delta>0$, independientes
de $U$, tales que, para todo $U\ge1$, se puede elegir un entero $k\ge U$
con
\begin{equation}
 \int_{[0,\infty)\setminus W_U}
 \frac{u^ke^{-(\sigma-2)u}}{k!}\dd u
 \le C_\delta\sigma^{-k}.
 \label{vent-cota-colas}
\end{equation}
Además, la elección satisface $k\asymp_\delta U$.
\end{lema}

\begin{proof}
Para $s>\ell>0$ y $0<U\le V$, separar una exponencial y aplicar
\eqref{vent-integral-gamma} proporciona
\begin{equation}
 \int_{[0,\infty)\setminus[U,V]}
 \frac{u^ke^{-su}}{k!}\dd u
 \le e^{-\ell V}(s-\ell)^{-k-1}
       +e^{\ell U}(s+\ell)^{-k-1}.
 \label{vent-colas-elementales}
\end{equation}
En la cola superior se usa $e^{-\ell u}\le e^{-\ell V}$; en la inferior,
$e^{\ell u}\le e^{\ell U}$. Las integrales restantes se amplían a toda
la semirrecta positiva.

Tomemos las constantes explícitas
\[
 \ell=2+\frac6\delta,\qquad b=1+\frac\delta2,\qquad
 M=\ell(1+\delta),\qquad \sigma=M(\ell+2),
\]
y el entero $k=\lceil b\sigma U\rceil$. Entonces $\sigma-2>\ell$ y
\begin{equation}
 b(\ell+2)=M-1,\qquad b(\ell-2)=\ell+1,\qquad
 b\sigma=M(M-1).
 \label{vent-identidades}
\end{equation}
La definición del techo implica
\begin{equation}
 \begin{gathered}
 b\sigma U\le k\le b\sigma U+1,\\
 U\le k\le(b\sigma+1)U\qquad(U\ge1).
 \end{gathered}
 \label{vent-orden-k}
\end{equation}

Pongamos $a_-=\sigma-2-\ell>0$ y $a_+=\sigma-2+\ell>0$.
La desigualdad $\log(\sigma/a)\le(\sigma-a)/a$ muestra que, para $a>0$,
\begin{equation}
 x+k\frac{\sigma-a}{a}\le c
 \quad\Longrightarrow\quad
 e^xa^{-k-1}\le\frac{e^c}{a\sigma^k}.
 \label{vent-comparacion-exponencial}
\end{equation}
Para el primer término de \eqref{vent-colas-elementales}, con
$s=\sigma-2$ y $V=(1+\delta)U$, las identidades
\eqref{vent-identidades} dan
\[
 -\ell(1+\delta)U+k\frac{\sigma-a_-}{a_-}
 =-MU+\frac{k}{M-1}\le\frac1{M-1}\le1.
\]
Para el segundo, observamos que $\ell a_+\le(\ell+1)\sigma$, pues esta
desigualdad equivale a $\ell(\ell-2)\le\sigma$, que se sigue de
$\sigma=\ell(1+\delta)(\ell+2)$. Así,
\[
 \ell U+k\frac{\sigma-a_+}{a_+}
 =\ell U-\frac{k(\ell-2)}{a_+}
 \le\ell U-\frac{(\ell+1)\sigma U}{a_+}\le0.
\]
Aplicar \eqref{vent-comparacion-exponencial} a ambas colas prueba
\eqref{vent-cota-colas}, con
\begin{equation}
 C_\delta=\frac e{\sigma-2-\ell}+\frac1{\sigma-2+\ell}.
 \label{vent-constante-colas}
\end{equation}
Todas las constantes se fijaron a partir de $\delta$, antes de conocer
la ventana. Esto prueba el lema.
\end{proof}

El núcleo normalizado
$(\sigma-2)^{k+1}u^ke^{-(\sigma-2)u}/k!$ tiene integral $1$.
Por \eqref{vent-cota-colas}, su masa exterior está acotada por
\[
 C_\delta(\sigma-2)
 \left(\frac{\sigma-2}{\sigma}\right)^k.
\]
La concentración es exponencial en $k$. Esa tasa, y no solo que la masa
exterior tienda a cero, permitirá comparar el núcleo con la contribución
de un polo.

{\raggedright
Véanse A, Lema~2.3, pp.~4--5, y \texttt{landau\_tails}.
B, Lema~2.2, p.~16, obtiene el mismo tipo de control mediante momentos
exponenciales de una densidad gamma.
\par}

\subsection{La cota superior para las derivadas}

Fijemos $\sigma,C_\delta$ como en el Lema~\ref{vent-colas}
y escribamos
\[
 E=\sigma-\theta,\qquad 0<E<\sigma.
\]
Para $v\ge0$, la serie exponencial da $v^k/k!\le e^v$.
Tomando $v=Eu$, resulta
\begin{equation}
 \frac{u^ke^{-\sigma u}}{k!}\le E^{-k}e^{-\theta u}.
 \label{vent-maximo-nucleo}
\end{equation}
En una ventana admisible, $g_+\le R_U$; fuera de ella,
$g_+\le C_0e^{2u}$. Aplicamos \eqref{vent-maximo-nucleo},
\eqref{vent-hipotesis} y el lema de colas a \eqref{vent-unilateral}.
Para el entero $k=\lceil b\sigma U\rceil$ obtenemos
\begin{equation}
 |A_k(\sigma+it)|+A_k(\sigma)
 \le 2e^{\eta U}E^{-k}+2C_0C_\delta\sigma^{-k}
 \qquad(t\in\R).
 \label{vent-cota-superior}
\end{equation}
La misma cota vale simultáneamente para todos los centros de la recta
$\Ree z=\sigma$. No hemos fijado todavía $\eta$: lo elegiremos después
de localizar el polo que contradiga esta estimación.

{\raggedright
Véanse A, (2.11), p.~5; \texttt{landau\_kernel\_max},
\texttt{landau\_window\_ineq} y \texttt{landau\_window\_key}.
\par}

\subsection{Elegir un polo mediante discos compactos}

Supongamos que existe $\rho_0\in Z$ con $\Ree\rho_0>\theta$.
No suponemos que haya un polo cuya parte real sea máxima.

El disco cerrado $\overline B(\sigma,E)$ queda dentro de $\Ree s>0$,
porque su extremo izquierdo tiene parte real $\theta>0$.
Su intersección $F$ con $Z$ es finita. Si $F$ es vacía, ponemos $D=E$.
Si no, ponemos
\[
 D=\min_{\rho\in F}|\sigma-\rho|.
\]
En ambos casos $0<D\le E$ y $B(\sigma,D)$ no contiene polos.
Consideramos dos alternativas.

\paragraph{Alternativa I.}
Existe un polo $\rho_1$ tal que $\sigma-\Ree\rho_1<D$.

\paragraph{Alternativa II.}
No existe tal polo. Esta posibilidad requiere $F\ne\varnothing$, pues,
si $D=E$ y $F$ fuera vacía, el polo $\rho_0$ satisfaría la alternativa I.
Un elemento que alcance la distancia mínima verifica
\[
 D=|\sigma-\rho|\ge\sigma-\Ree\rho\ge D.
\]
La igualdad fuerza $\rho=\beta:=\sigma-D\in\R$.
El punto más cercano es único. Además,
\[
 D\le\sigma-\Ree\rho_0<E,
\]
de modo que $\beta>\theta$. Como hay finitos polos en el disco cerrado,
existe $L_\beta$ con $D<L_\beta\le E$ para el cual $\beta$ es el único
polo de $B(\sigma,L_\beta)$.

Esto demuestra la dicotomía. En la segunda alternativa ya tenemos un
polo real aislado; falta precisar la elección de centro en la primera.

\subsubsection*{Un desplazamiento vertical elimina los empates}

En la alternativa I, sean $h=\sigma-\Ree\rho_1>0$ y
$t_0=\operatorname{Im}\rho_1$. Elegimos $\varepsilon>0$ tan pequeño que
\[
 h+\varepsilon<D,\qquad D+\varepsilon<\sigma.
\]
El conjunto $F_1=Z\cap\overline B(\sigma+it_0,D+\varepsilon)$ es finito.
Si $|t-t_0|<\varepsilon$ y $z=\sigma+it$, entonces $|z-\rho_1|<D$,
mientras que cualquier polo fuera de $F_1$ está a distancia mayor que
$D$ de $z$. Por tanto, el polo más cercano debe pertenecer a $F_1$.

Para dos polos distintos $a,b\in F_1$, la diferencia de cuadrados de sus
distancias a $\sigma+it$ es una función afín de $t$, con pendiente
$-2(\operatorname{Im}a-\operatorname{Im}b)$.
Si sus partes imaginarias son distintas, hay a lo sumo un valor de
empate. Si coinciden, el empate exigiría
\[
 (\sigma-\Ree a)^2=(\sigma-\Ree b)^2.
\]
Ambas distancias horizontales son positivas, por lo que tendríamos
$a=b$, contradicción. Podemos evitar los finitos valores de empate
dentro de $(t_0-\varepsilon,t_0+\varepsilon)$.

Fijamos uno de esos $t$. Existe entonces un único polo más cercano
$\rho$, con
\[
 0<d:=|\sigma+it-\rho|<D\le E.
\]
La finitud del conjunto relevante permite elegir un radio $L$ con
$d<L\le D$ tal que $B(\sigma+it,L)$ contenga solo ese polo.
Tomamos un radio interior $r$ con $d<r<L$. Así, $\Phi$ es holomorfa
cerca del disco cerrado $\overline B(\sigma,r)$, y, tras quitar la
parte principal en $\rho$, también lo será cerca de
$\overline B(\sigma+it,r)$.

En esta alternativa el polo elegido es necesariamente no real: si fuera
real, su distancia a $\sigma+it$ sería al menos su distancia a $\sigma$,
que es mayor o igual que $D$, en contradicción con $d<D$.
Para la estimación posterior bastará que su residuo sea no nulo.
En toda la selección se minimizó sobre conjuntos finitos dentro de compactos.

{\raggedright
Véanse A, Lema~2.4 y demostración del Teorema~2.2, pp.~5--6;
B, pp.~17--18; \texttt{landau\_finite}, \texttt{landau\_dichotomy}
y \texttt{landau\_nearest}.
\par}

\subsection{Cauchy hace visible el polo}

En la alternativa I, escribamos $z=\sigma+it$ y sea $c\ne0$
el residuo de $\Phi$ en $\rho$. La función
\[
 h_\rho(w)=\Phi(w)-\frac c{w-\rho}
\]
se prolonga holomórficamente en $\rho$ y en todo $B(z,L)$.
Aplicando Cauchy en la circunferencia de radio interior $r<L$, obtenemos
constantes $M_z,M_\sigma\ge0$ independientes de $k$ tales que
\begin{equation}
 \left|A_k(z)-\frac c{(z-\rho)^{k+1}}\right|\le M_zr^{-k},
 \qquad |A_k(\sigma)|\le M_\sigma r^{-k}.
 \label{vent-cauchy}
\end{equation}
Por ejemplo, $M_z$ es el máximo de $|h_\rho|$ en $|w-z|=r$;
ese máximo es finito por continuidad y compacidad.
El factor del polo procede de
\[
 \frac{\mathrm d^k}{\mathrm dw^k}\frac c{w-\rho}
 =\frac{(-1)^kk!c}{(w-\rho)^{k+1}}.
\]
La desigualdad triangular inversa y la realidad de $A_k(\sigma)$ dan
\begin{equation}
 |A_k(z)|+A_k(\sigma)
 \ge\frac{|c|}{d^{k+1}}-(M_z+M_\sigma)r^{-k}.
 \label{vent-cota-polo}
\end{equation}
El valor absoluto elimina la fase compleja de la contribución dominante;
basta que $c\ne0$.

En la alternativa II ponemos $d=\sigma-\beta$ y elegimos
$d<r<L_\beta$. El residuo $c_\beta$ es real positivo.
Al quitar la parte principal y aplicar Cauchy se obtiene
\begin{equation}
 A_k(\sigma)\ge\frac{c_\beta}{d^{k+1}}-Mr^{-k}.
 \label{vent-polo-real}
\end{equation}
La expresión de la izquierda de \eqref{vent-cota-superior}, ahora
con $t=0$, es al menos $A_k(\sigma)$. Por tanto, en ambas alternativas
existen constantes $a>0$, $M\ge0$ y distancias $0<d<r<E$ tales que
\begin{equation}
 |A_k(\sigma+it)|+A_k(\sigma)
 \ge a\,d^{-k}-Mr^{-k}
 \qquad(k\ge0).
 \label{vent-cota-inferior}
\end{equation}
En la primera alternativa $a=|c|/d$; en la segunda $a=c_\beta/d$.
Los centros, radios y constantes ya están fijados antes de elegir
una ventana. El radio interior garantiza que Cauchy se aplica sobre
una circunferencia contenida en la región de holomorfía.

{\raggedright
Véanse A, p.~6; B, pp.~17--18; \texttt{landau\_cauchy}
y \texttt{landau\_pole\_asym}.
\par}

\subsection{Contradicción y papel del residuo positivo}

\begin{proof}[Conclusión de la demostración de la Proposición~\ref{vent-criterio}]
Fijemos ahora
\[
 \eta=\frac12\log(E/d)>0.
\]
La hipótesis de ventanas proporciona, para este $\eta$, valores de $U$
arbitrariamente grandes y los correspondientes $k\ge U$.
Al combinar \eqref{vent-cota-superior} y \eqref{vent-cota-inferior}
y multiplicar por $d^k$, resulta
\begin{equation}
 a\le M\left(\frac dr\right)^k
 +2e^{\eta U}\left(\frac dE\right)^k
 +2C_0C_\delta\left(\frac d\sigma\right)^k.
 \label{vent-contradiccion}
\end{equation}
Como $U\le k$,
\[
 e^{\eta U}\left(\frac dE\right)^k
 \le\left(e^\eta\frac dE\right)^k
 =\left(\frac dE\right)^{k/2}.
\]
Las tres bases exponenciales son estrictamente menores que $1$.
Al tomar ventanas con $U\to\infty$, también $k\to\infty$, y el lado
derecho de \eqref{vent-contradiccion} tiende a cero.
Esto contradice $a>0$ y demuestra la proposición.
\end{proof}

El orden de elección fue: datos $g,\theta,\delta$; constantes del núcleo;
un polo hipotético, centros, radios y constantes de Cauchy; el margen
$\eta$; finalmente una ventana suficientemente lejana y su orden $k$.
No hace falta controlar los huecos entre ventanas ni disponer de la
estimación para todos los órdenes de derivación.
Véanse A, p.~6; \texttt{landau\_window} y
\texttt{landau\_growth\_contra}.

La hipótesis de signo sobre los polos reales es necesaria.
Si $0<\theta<\beta<1$ y $g(u)=-e^{\beta u}$, entonces $g(u)\le0$
en toda ventana, pero
\[
 \Phi(s)=-\frac1{s-\beta}
\]
tiene un polo en $\beta>\theta$ con residuo $-1$.
En ese ejemplo, $A_k(\sigma)<0$ y
$|A_k(\sigma)|+A_k(\sigma)=0$. El residuo positivo de la
alternativa II impide precisamente esta cancelación.

\subsection{La versión puntual y su aplicación a las funciones L}

\begin{proposicion}[Corolario: cota puntual]
\label{vent-corolario}
Supongamos las hipótesis analíticas de la Proposición~\ref{vent-criterio}.
Si hay una constante fija $C\in\R$ y valores arbitrariamente grandes
de $Y$ tales que
\begin{equation}
 H(x)\le Cx^\theta
 \qquad\text{para todo }x\in[Y,Y^{1+\delta}],
 \label{vent-cota-puntual}
\end{equation}
donde $g(u)=H(e^u)$, entonces $\Phi$ no tiene polos con parte real
mayor que $\theta$.
\end{proposicion}

\begin{proof}
Reemplazamos $C$ por $C_+=\max\{C,0\}$ y ponemos $U=\log Y$.
La ventana se transforma en $[U,(1+\delta)U]$ y
$g(u)\le C_+e^{\theta u}$. Tomando $R_U(u)=C_+e^{\theta u}$,
\[
 \int_U^{(1+\delta)U}R_U(u)e^{-\theta u}\dd u=C_+\delta U.
\]
Para cada $\eta>0$, esta cantidad es a lo sumo $e^{\eta U}$ cuando
$U$ es suficientemente grande. Por ejemplo, basta
$U\ge\max\{1,2C_+\delta/\eta^2\}$, pues
$e^{\eta U}\ge\eta^2U^2/2$.
Como las ventanas son arbitrariamente lejanas, se puede imponer
ese umbral. Aplicamos la Proposición~\ref{vent-criterio}.
\end{proof}

Hay también una ruta directa, utilizada por B: dentro de la ventana,
\eqref{vent-integral-gamma} proporciona
\begin{equation}
 \begin{aligned}
 |A_k(\sigma+it)|+A_k(\sigma)
 &\le 2C_+E^{-k-1}+2C_0C_\delta\sigma^{-k}\\
 &\le\left(\frac{2C_+}{E}+2C_0C_\delta\right)E^{-k}.
 \end{aligned}
 \label{vent-ruta-puntual}
\end{equation}
La comparación con \eqref{vent-cota-inferior} conduce a la misma
contradicción sin introducir $\eta$. La versión integrada permite
cotejar el enunciado formal más general.
Véanse A, Corolario~2.5, p.~6, y B, Proposición~2.1 y (2.4),
pp.~14--18.

Para una familia finita de caracteres primitivos y pesos $a_\chi>0$,
la construcción de la Nota~\ref{sec:sumas} da
\[
 H(x)=\sum_\chi a_\chi\Ree G_\chi(x),\qquad
 \Phi(s)=\frac12\sum_\chi a_\chi
 \bigl(F_\chi(s)+F_{\bar\chi}(s)\bigr).
\]
El crecimiento $|g(u)|\le C_0e^{2u}$ procede de la cota elemental de las
sumas suavizadas. La continuación meromorfa y la regularidad en $1$
proceden de sus fórmulas de Mellin. En un cero representado por la
familia, el residuo es
\begin{equation}
 \begin{gathered}
 \Res_{s=\rho}\Phi(s)=\frac{b_\rho}{\rho(1-\rho)},\\
 b_\rho=\frac12\sum_\chi a_\chi
 \bigl(m_\chi(\rho)+m_{\bar\chi}(\rho)\bigr)>0.
 \end{gathered}
 \label{vent-residuo-familia}
\end{equation}
Los pesos positivos evitan la cancelación de polos. En un polo real
$\beta\in(0,1)$, también dan el signo positivo del residuo.
Para situar los polos reales en ese intervalo bastan el producto
de Euler y la regularidad en $1$; no se supone previamente la
hipótesis de Riemann.

Aplicamos el corolario a $H_\psi$ y $\delta=1/2$.
Si existe $\varepsilon_0>0$ tal que la construcción aritmética entrega
\eqref{vent-cota-puntual} con $\theta=1/2+\varepsilon$ para cada
$0<\varepsilon<\varepsilon_0$, cualquier cero con parte real
$\beta>1/2$ queda excluido al elegir
\[
 0<\varepsilon<\min\{\varepsilon_0,\,\beta-1/2,\,1/2\}.
\]
La ecuación funcional permite tratar después el lado izquierdo de la
franja. Para caracteres no primitivos se retiran los factores de Euler
adicionales, no nulos en $\Ree s>0$. Estas últimas reducciones se
identifican en A, pp.~6--7, y B, Proposición~2.3, pp.~18--19.
La construcción aritmética de las ventanas sigue pendiente de desarrollo
en el cuaderno: haber demostrado el criterio no suministra esa estimación.

\subsection{Concordancia con Lean y alcance del cotejo}

Se consultó el espejo
\href{https://github.com/CocoSgt/riemann-hypothesis-lean/tree/7cbf7a0c32fca1badf05931879a37c3f8a2b03b4}
{\texttt{CocoSgt/riemann-hypothesis-lean}},
en la revisión
\href{https://github.com/CocoSgt/riemann-hypothesis-lean/tree/7cbf7a0c32fca1badf05931879a37c3f8a2b03b4}
{\texttt{7cbf7a0c32fc}},
que declara Lean \texttt{4.33.1} y Mathlib
\texttt{0df444a360ea}. Los identificadores completos se conservan en
el registro del repositorio. Esta referencia fija el código leído;
no afirma que hayamos comparado todo el espejo con el archivo original
distribuido por Alpöge.

\begingroup
\raggedright
\begin{description}
 \item[Datos analíticos.]
 \texttt{PoleData} incluye medibilidad, crecimiento, identidad de Laplace,
 holomorfía fuera de $Z$, ausencia de acumulación local y una parte
 principal simple con residuo real positivo cuando el polo es real.

 \item[Hipótesis integrada.]
 \texttt{WindowHypU} exige la integrabilidad del mayorante ponderado de
 \eqref{vent-hipotesis} y ventanas arbitrariamente lejanas para cada
 $\eta>0$.

 \item[Derivadas y núcleo.]
 \texttt{iteratedDeriv\_laplace} y \texttt{kernel\_integral} identifican
 las derivadas y la integral del núcleo. La representación integral
 de las derivadas se utiliza solo en $\Ree z>2$.

 \item[Control uniforme.]
 \texttt{landau\_tails}, \texttt{landau\_kernel\_max} y
 \texttt{landau\_window\_ineq} fijan las constantes antes de la ventana
 y conservan el signo del término real.

 \item[Geometría.]
 \texttt{landau\_pole\_re\_le}, \texttt{landau\_finite},
 \texttt{landau\_dichotomy} y \texttt{landau\_nearest} obtienen
 $\Ree\rho\le2$ y efectúan la selección en conjuntos finitos.

 \item[Estimaciones locales.]
 \texttt{landau\_cauchy} y \texttt{landau\_pole\_asym} restan la parte
 principal y aplican Cauchy con un radio estrictamente interior.

 \item[Conclusión.]
 \texttt{landau\_window\_key}, \texttt{landau\_window\_contra} y
 \texttt{landau\_window} comparan la cota de derivadas con la contribución
 de un polo de parte real mayor que $\theta$.

 \item[Puente desde cotas puntuales.]
 \texttt{windowHyp\_of\_rpow}, \texttt{windowHypU\_of\_windowHyp}
 y \texttt{zeroFree\_of\_rpow\_window}
 deducen la condición integrada de la cota puntual y la transportan
 a \texttt{WindowHypU} mediante el cambio $x=e^u$; el corolario
 concluye la no anulación.
\end{description}
\par
\endgroup

{\raggedright
Las definiciones están en
\href{https://github.com/CocoSgt/riemann-hypothesis-lean/blob/7cbf7a0c32fca1badf05931879a37c3f8a2b03b4/src/Development/AutomorphicArithmetic/DedekindLFunctions/LandauZeroFreeCriterion.lean}
{\texttt{LandauZeroFreeCriterion.lean}};
el ensamble final está en
\href{https://github.com/CocoSgt/riemann-hypothesis-lean/blob/7cbf7a0c32fca1badf05931879a37c3f8a2b03b4/src/Development/AutomorphicArithmetic/DedekindLFunctions/LandauTheorems/landauWindow.lean}
{\texttt{landauWindow.lean}}.
La definición llamada
\texttt{Deduction.}\allowbreak\texttt{LandauZeroFreeCriterion},
al final del primer archivo, expresa otra variante con una cota eventual
sobre una semirrecta; el teorema de ventanas estudiado aquí es
\texttt{Deduction.landau\_window}.
El registro del repositorio conserva las rutas, declaraciones y hashes
de los archivos recuperados.
\par}

El criterio general no incluye como hipótesis la ausencia de polos en
$\Ree s>1$ ni la regularidad en $1$; estas se satisfacen en la aplicación
a las funciones $L$. \texttt{WindowHypU} exige un mayorante puntual
no negativo de $g$ y controla su tamaño mediante la integral ponderada;
no exige una cota puntual de la forma $g(u)\le Ce^{\theta u}$.
A, B y Lean concuerdan en el mecanismo revisado, con esas diferencias
de formulación.

\paragraph{Estado del tramo.}
Se ha reconstruido la demostración matemática y se ha efectuado un
cotejo localizado con los enunciados y pruebas señalados.
Este cotejo no incluyó compilación ni ejecución de Lean, ni una auditoría
de todas sus dependencias. Esta nota cierra el pendiente analítico del
cuaderno; el siguiente desarrollo es la identidad de caracteres y
la cancelación de su término principal.

\clearpage
\section{Identidad de caracteres y cancelación del término principal}
\label{sec:caracteres}

Esta nota reconstruye el enlace entre una condición de Frobenius y las sumas corregidas de la Nota~\ref{sec:sumas}. La identidad principal expresa una suma aritmética no negativa, después de sustraer su término principal, como combinación de las funciones $G_\chi$. Después estudiamos los promedios finitos que seleccionan caracteres y explican el coeficiente positivo de $H_\psi$.

\fuente{A, §4, especialmente Proposición~4.2 y ecuaciones (4.5)--(4.6), pp.~13--17; B, §3, especialmente Proposición~3.1, pp.~19--23. Las barras de conjugación y los signos se comprobaron también en las páginas renderizadas de ambos PDF. El cotejo con la revisión fijada de Lean se identifica al final.}

La parte finita se demuestra aquí, incluido el fundamento de la ortogonalidad. La interpretación mediante extensiones abelianas y caracteres primitivos se señala como entrada clásica. La comparación con periodos automorfos y la producción de las ventanas permanecen como desarrollos posteriores. El resultado analítico que utilizará esas ventanas está en la Nota~\ref{sec:ventanas}.

\subsection{Caracteres de un grupo abeliano finito}

Sea $\Gamma$ un grupo abeliano finito, escrito aditivamente. Un carácter es un homomorfismo $\chi:\Gamma\to\C^\times$. Para todo $g\in\Gamma$, si $ng=0$, entonces $\chi(g)^n=1$. Por tanto, sus valores tienen módulo uno y

\[
\chi(-g)=\chi(g)^{-1}=\overline{\chi(g)}.
\]

El conjunto $\widehat\Gamma$ de todos los caracteres forma un grupo con el producto punto a punto. Su elemento identidad es el carácter constante $1$.

\subsubsection*{Extender un carácter y contar todos los caracteres}

Se puede justificar la ortogonalidad sin presuponer la dualidad de grupos finitos. Sea $H\le\Gamma$, sea $\psi:H\to\C^\times$ un carácter y sea $b\in\Gamma\setminus H$. Denotemos por $n$ el orden de $b+H$ en $\Gamma/H$. Así, $n\ge2$, $nb\in H$ y

\[
H'=H+\Z b,\qquad [H':H]=n.
\]

Elijamos una raíz $z$ de $z^n=\psi(nb)$ y definamos

\begin{equation}
\widetilde\psi(h+jb)=\psi(h)z^j,
\qquad h\in H,\ j\in\Z.
\label{car-extension}
\end{equation}

Para comprobar que la definición es independiente de los representantes, supongamos que $h+jb=h'+j'b$. Entonces $(j-j')b=h'-h\in H$, de modo que $j-j'=kn$ para algún $k\in\Z$. En consecuencia, $h'=h+knb$ y

\[
\psi(h')z^{j'}
=\psi(h)\psi(nb)^kz^{j'}
=\psi(h)z^{kn+j'}
=\psi(h)z^j.
\]

La fórmula \eqref{car-extension} respeta la suma y restringe a $\psi$ en $H$, por lo que define un carácter de $H'$. Toda extensión tiene esta forma: su valor en $b$ determina los restantes valores y debe satisfacer $z^n=\psi(nb)$. Como $\psi(nb)\ne0$, esta ecuación tiene exactamente $n$ raíces complejas distintas. Por tanto, \textbf{cada carácter de $H$ tiene exactamente $[H':H]$ extensiones a $H'$}.

Comenzando en $H=\{0\}$ y adjuntando sucesivamente elementos hasta obtener $\Gamma$, el número de caracteres se multiplica en cada paso por el índice del nuevo subgrupo. El producto de estos índices es $|\Gamma|$, de donde

\begin{equation}
|\widehat\Gamma|=|\Gamma|.
\label{car-cardinal-dual}
\end{equation}

Este mismo procedimiento demuestra que los caracteres separan puntos. Si $v\ne0$ tiene orden $d\ge2$, el carácter de $\langle v\rangle$ definido por

\[
jv\longmapsto e^{2\pi i j/d}
\]

se extiende a $\Gamma$ y toma en $v$ un valor distinto de uno. Solo se han utilizado aquí la existencia y el número de las raíces complejas de un número no nulo, junto con propiedades elementales de los grupos finitos.

\subsubsection*{Las dos formas de ortogonalidad}

Si $H\le\Gamma$ y $\chi\in\widehat\Gamma$, entonces

\begin{equation}
\frac1{|H|}\sum_{h\in H}\chi(h)
=
\begin{cases}
1,&\chi|_H=1,\\
0,&\chi|_H\ne1.
\end{cases}
\label{car-ortogonalidad-subgrupo}
\end{equation}

La primera posibilidad es inmediata. En la segunda, elijamos $h_0\in H$ con $\chi(h_0)\ne1$. La traslación $h\mapsto h+h_0$ permuta $H$; por eso la suma es igual a sí misma multiplicada por $\chi(h_0)$ y debe ser cero. La misma fórmula vale con $\overline{\chi(h)}$.

Por otra parte, para todo $v\in\Gamma$,

\begin{equation}
\sum_{\chi\in\widehat\Gamma}\chi(v)
=|\Gamma|\,\mathbf 1_{\{v=0\}}.
\label{car-ortogonalidad-dual}
\end{equation}

Si $v=0$, cada sumando es uno y se aplica \eqref{car-cardinal-dual}. Si $v\ne0$, existe $\psi\in\widehat\Gamma$ con $\psi(v)\ne1$. La permutación $\chi\mapsto\psi\chi$ de $\widehat\Gamma$ multiplica la suma por $\psi(v)$, de modo que vuelve a ser cero. En particular,

\begin{equation}
\sum_{\chi\in\widehat\Gamma}
\overline{\chi(g)}\chi(v)
=|\Gamma|\,\mathbf1_{\{v=g\}}.
\label{car-ortogonalidad-pares}
\end{equation}

El conjugado convierte el producto del lado izquierdo en $\chi(v-g)$. Si se omite, el producto es $\chi(v+g)$ y se selecciona $v=-g$. Para caracteres no reales, estas condiciones pueden ser distintas.

\subsection{Inercia, Frobenius y los coeficientes primitivos}

Seguimos primero con datos finitos abstractos. Para cada primo racional $p$, fijamos un subgrupo $I_p\le\Gamma$ y una clase $F_p\in\Gamma/I_p$. Si $\widetilde F_p\in\Gamma$ es un representante, definimos, para cada entero $j\ge1$,

\begin{equation}
a_\chi(p^j)
:=\frac1{|I_p|}\sum_{h\in I_p}
\chi(j\widetilde F_p+h).
\label{car-media-inercia}
\end{equation}

Cambiar $\widetilde F_p$ por $\widetilde F_p+h_0$ con $h_0\in I_p$ cambia $h$ por $h+jh_0$. Como esta traslación permuta $I_p$, la media es independiente del representante. Además,

\begin{equation}
a_\chi(p^j)
=\chi(\widetilde F_p)^j
\frac1{|I_p|}\sum_{h\in I_p}\chi(h)
=
\begin{cases}
\chi(F_p)^j,&\chi|_{I_p}=1,\\
0,&\chi|_{I_p}\ne1.
\end{cases}
\label{car-coeficiente-local}
\end{equation}

En la primera línea, $\chi(F_p)$ significa el valor del carácter descendido a $\Gamma/I_p$; ese valor no requiere elegir un representante. Extendemos $a_\chi$ multiplicativamente a los enteros positivos, con $a_\chi(1)=1$. En particular, $a_1(n)=1$ para todo $n\ge1$, incluso en los primos donde $I_p$ sea no trivial. Estas afirmaciones y la independencia del representante quedan demostradas solo con \eqref{car-ortogonalidad-subgrupo}. Corresponden a B, §3.1, pp. 19–20.

\subsubsection*{Interpretación aritmética: entrada externa}

Cuando $E/\Q$ es una extensión finita abeliana, se toma $\Gamma=\operatorname{Gal}(E/\Q)$. Para un primo $\mathfrak p$ de $E$ sobre $p$, su grupo de descomposición $D_{\mathfrak p}$ actúa en el cuerpo residual. La inercia $I_{\mathfrak p}$ es el núcleo de esta acción. El cociente $D_{\mathfrak p}/I_{\mathfrak p}$ es el grupo de Galois de la extensión residual y contiene el Frobenius aritmético, que actúa como $x\mapsto x^p$. Su clase determina $F_p\in\Gamma/I_p$. La conmutatividad de $\Gamma$ hace independientes de $\mathfrak p$ tanto $I_p$ como esta clase de Frobenius.

Usamos la convención de Frobenius aritmético para que en $\Q(\zeta_N)$, cuando $p\nmid N$, $F_p$ corresponda a $\zeta_N\mapsto\zeta_N^p$. Con esta convención se obtiene el valor usual de un carácter de Dirichlet en $p$. La convención inversa cambia también los valores que deben aparecer en las fórmulas.

La identificación aritmética utilizada en B es que a cada $\chi\in\widehat\Gamma$ le corresponde un carácter de Dirichlet \textbf{primitivo} $\chi^*$, de conductor $q_\chi$, tal que

\begin{equation}
a_\chi(n)=\chi^*(n),\qquad n\ge1.
\label{car-caracter-primitivo}
\end{equation}

Los ingredientes clásicos de \eqref{car-caracter-primitivo} son la realización ciclotómica de las extensiones abelianas de $\Q$ y la descripción de sus caracteres y factores de Euler locales. Localmente, la afirmación es

\begin{equation}
\chi|_{I_p}=1\ \Longleftrightarrow\ p\nmid q_\chi;
\qquad
\chi(F_p)=\chi^*(p)\quad\text{si }p\nmid q_\chi.
\label{car-conductor}
\end{equation}

Si $\chi$ no es trivial sobre la inercia, $p\mid q_\chi$, y ambos lados de \eqref{car-caracter-primitivo} se anulan en $p^j$. Si sí lo es, \eqref{car-coeficiente-local} y \eqref{car-conductor} identifican los valores en todas las potencias de $p$. La multiplicatividad concluye la igualdad para $n\ge1$. La deducción de \eqref{car-caracter-primitivo} a partir de \eqref{car-conductor} es elemental; \eqref{car-conductor} es la entrada de teoría de números que se cita, sin reconstruir aquí la teoría de las extensiones abelianas.

Esta precisión importa en un primo ramificado en $E$: puede ocurrir que un carácter concreto sea trivial en $I_p$. Su factor local y su coeficiente $a_\chi(p)$ permanecen entonces presentes. La ramificación del campo común no obliga a que todos sus caracteres se anulen en ese primo.

\subsubsection*{Un carácter inducido a un módulo mayor}

Si un carácter primitivo $\chi^*$ de conductor $q$ induce un carácter $\chi_N$ módulo $N$, con $q\mid N$, entonces

\[
\chi_N(n)=
\begin{cases}
\chi^*(n),&\gcd(n,N)=1,\\
0,&\gcd(n,N)>1.
\end{cases}
\]

En un primo $p\mid N$ con $p\nmid q$, el carácter inducido vale cero, aunque $\chi^*(p)$ tiene módulo uno. Por los productos de Euler, inicialmente en $\Ree s>1$,

\begin{equation}
L(s,\chi_N)
=L(s,\chi^*)
\prod_{\substack{p\mid N\\p\nmid q}}
(1-\chi^*(p)p^{-s}).
\label{car-euler-inducido}
\end{equation}

La fórmula se encuentra también al final de B, §2, p. 19. La media de inercia \eqref{car-media-inercia} produce $\chi^*$, por lo que no se puede sustituir $a_\chi(n)$ por el carácter inducido a un módulo común sin llevar la cuenta de estos factores locales. El carácter trivial de $\widehat\Gamma$ corresponde al carácter primitivo de conductor uno y a $\zeta(s)$.

\subsection{Del promedio de inercia a una condición de Frobenius}

Sea $\kappa:\Gamma\to\{\pm1\}$ un carácter cuadrático no trivial y pongamos

\[
\Phi=\ker\kappa,\qquad f=|\Phi|,
\qquad |\Gamma|=2f.
\]

Para cada $g\in\Gamma$, $p$ primo y $j\ge1$, aplicamos \eqref{car-ortogonalidad-pares} a cada término de \eqref{car-media-inercia}:

\begin{equation}
\begin{aligned}
\frac1f\sum_{\chi\in\widehat\Gamma}
\overline{\chi(g)}\,a_\chi(p^j)
&=\frac1{f|I_p|}\sum_{h\in I_p}
\sum_{\chi\in\widehat\Gamma}
\chi(j\widetilde F_p+h-g)\\
&=\frac{|\Gamma|}{f|I_p|}
\#\{h\in I_p:j\widetilde F_p+h=g\}\\
&=\frac2{|I_p|}\,
\mathbf1_{\{g-j\widetilde F_p\in I_p\}}.
\end{aligned}
\label{car-frobenius}
\end{equation}

La ecuación para $h$ tiene una única solución cuando $g-j\widetilde F_p\in I_p$, y no tiene ninguna en caso contrario. La condición es independiente del representante; equivale a $g+I_p=jF_p$ en $\Gamma/I_p$. Esta es exactamente B, ecuación (3.1), p. 20.

El factor $2$ procede de $|\Gamma|/f=2$. El denominador $|I_p|$ procede de la media sobre la inercia. El resultado es real y no negativo, aunque cada término de la suma de caracteres pueda ser complejo.

\subsection{La suma ponderada y su positividad}

Fijemos $c\in\Gamma\setminus\Phi$ y $t\in\Phi$. La aplicación

\[
\rho_t:\Gamma\setminus\Phi\longrightarrow\Phi,
\qquad \rho_t(g)=c+g+t
\]

es una biyección: es una traslación y $\kappa(c+g+t)=(-1)(-1)(1)=1$. Su inversa en $\Phi$ es $\rho\mapsto\rho-c-t$. No se necesita suponer que $c$ tenga orden dos para esta identidad finita.

Sean $a>0$ y $w_\rho>0$ para todo $\rho\in\Phi$. Escribimos $v_+=\max(v,0)$ para $v\in\R$ y definimos

\begin{equation}
\mathcal A(a,t,w)
:=\frac1f\sum_{g\notin\Phi}
\sum_{\chi\in\widehat\Gamma}\overline{\chi(g)}
\sum_{n\ge1}\Lambda(n)a_\chi(n)
\left(\frac an-w_{\rho_t(g)}\right)_+.
\label{car-suma-ponderada}
\end{equation}

Aquí $\Lambda(p^j)=\log p$ para $p$ primo y $j\ge1$, y $\Lambda(n)=0$ en los demás casos, incluido $n=1$. Es la ecuación (3.2) de B, p. 20. Como $\Phi$ es finito y los pesos son positivos, $w_{\min}:=\min_{\rho\in\Phi}w_\rho>0$. Todo término no nulo satisface $n<a/w_{\min}$. En el extremo $n=a/w_\rho$ la parte positiva vale cero. Por tanto, todas las sumas escritas son efectivamente finitas y se pueden reordenar sin hipótesis analíticas de convergencia.

Agrupando los $n$ con $\Lambda(n)\ne0$ y aplicando \eqref{car-frobenius}, se obtiene

\begin{equation}
\mathcal A(a,t,w)
=\sum_p(\log p)\sum_{j\ge1}
\frac2{|I_p|}
\sum_{\substack{g\notin\Phi\\g-j\widetilde F_p\in I_p}}
\left(\frac a{p^j}-w_{\rho_t(g)}\right)_+.
\label{car-suma-positiva}
\end{equation}

Todos los factores de cada sumando son reales y no negativos. En particular, $\mathcal A(a,t,w)\ge0$. La fórmula mantiene tanto la condición de Frobenius como el umbral $p^j<a/w_{\rho_t(g)}$ que interviene después en las estimaciones de divisibilidad de periodos.

\subsection{Términos corregidos y cancelación exacta}

Para cada $\chi\ne1$, fijemos por ahora una constante arbitraria $\ell_\chi\in\C$. Para \textbf{todo} $y>0$, definimos

\begin{equation}
\begin{aligned}
G_\chi(y)
&=\sum_{n\ge1}\Lambda(n)a_\chi(n)(y/n-1)_+
+y\ell_\chi,\qquad \chi\ne1,\\
G_1(y)
&=\sum_{n\ge1}\Lambda(n)(y/n-1)_+
-y\log y+(1+\gamma_E)y.
\end{aligned}
\label{car-sumas-corregidas}
\end{equation}

La constante $\gamma_E$ es la constante de Euler que aparece en la regularización de $\zeta$. Las sumas finitas siguen teniendo sentido cuando $0<y<1$: en ese intervalo se anulan y solo quedan las correcciones indicadas. Esta extensión es necesaria porque en \eqref{car-suma-ponderada} no se exige $a/w_\rho\ge1$. Las definiciones son las de B, §3.2, p. 21; bajo \eqref{car-caracter-primitivo} y con $\ell_\chi=L'(1,\chi^*)/L(1,\chi^*)$, coinciden con las funciones de B, §1.2, p. 4.

\subsubsection*{El promedio en la clase donde $\kappa=-1$}

Para cada $\chi\in\widehat\Gamma$, se tiene

\begin{equation}
\frac1f\sum_{g\notin\Phi}\overline{\chi(g)}
=
\begin{cases}
1,&\chi=1,\\
-1,&\chi=\kappa,\\
0,&\chi\notin\{1,\kappa\}.
\end{cases}
\label{car-promedio-clase}
\end{equation}

En efecto, $\Gamma\setminus\Phi=c+\Phi$, por lo que el promedio es

\[
\overline{\chi(c)}\,
\frac1f\sum_{h\in\Phi}\overline{\chi(h)}.
\]

Por \eqref{car-ortogonalidad-subgrupo}, se anula salvo que $\chi|_\Phi=1$. En ese caso, $\chi$ desciende al grupo de dos elementos $\Gamma/\Phi$ y solo puede ser $1$ o $\kappa$. Sus valores en $c$ son respectivamente $1$ y $-1$, ambos reales. Esto conserva el conjugado y explica el signo negativo de la segunda línea.

\begin{minipage}{\linewidth}
\subsubsection*{Corrección de cada carácter antes de sumar}

Como $w>0$,

\[
\left(\frac an-w\right)_+
=w\left(\frac a{wn}-1\right)_+.
\]
\end{minipage}

Así, para $\chi\ne1$,

\begin{equation}
\sum_{n\ge1}\Lambda(n)a_\chi(n)
\left(\frac an-w\right)_+
=wG_\chi(a/w)-a\ell_\chi.
\label{car-correccion-no-principal}
\end{equation}

Para el carácter trivial, usando $a_1(n)=1$,

\begin{equation}
\sum_{n\ge1}\Lambda(n)
\left(\frac an-w\right)_+
=wG_1(a/w)+a\log a-a\log w-(1+\gamma_E)a.
\label{car-correccion-principal}
\end{equation}

El signo de $-a\ell_\chi$ en \eqref{car-correccion-no-principal} y los tres términos añadidos en \eqref{car-correccion-principal} se obtienen despejando las sumas a partir de \eqref{car-sumas-corregidas}.

\subsubsection*{La identidad principal: demostración completa}

Definamos

\begin{equation}
\mathcal M(a,w)
:=a\left(
\log a-1-\gamma_E+\ell_\kappa
-\frac1f\sum_{\rho\in\Phi}\log w_\rho
\right).
\label{car-termino-principal}
\end{equation}

Al sustituir \eqref{car-correccion-no-principal} y \eqref{car-correccion-principal} en \eqref{car-suma-ponderada}, los términos que contienen las funciones $G_\chi$ dan

\begin{equation}
\frac1f\sum_{g\notin\Phi}
\sum_{\chi\in\widehat\Gamma}
\overline{\chi(g)}\,w_{\rho_t(g)}
G_\chi\!\left(\frac a{w_{\rho_t(g)}}\right).
\label{car-resto}
\end{equation}

Las correcciones de los caracteres no triviales suman, por \eqref{car-promedio-clase},

\[
-a\sum_{\chi\ne1}\ell_\chi
\left(\frac1f\sum_{g\notin\Phi}\overline{\chi(g)}\right)
=-a\ell_\kappa(-1)
=a\ell_\kappa.
\]

Para el carácter trivial, la biyección $\rho_t$ da

\[
\begin{aligned}
\frac1f\sum_{g\notin\Phi}
\bigl[a\log a-a\log w_{\rho_t(g)}-(1+\gamma_E)a\bigr]
&=a\log a-(1+\gamma_E)a
-\frac af\sum_{\rho\in\Phi}\log w_\rho.
\end{aligned}
\]

La suma de ambas correcciones es exactamente \eqref{car-termino-principal}. Por tanto,

\begin{equation}
\boxed{
\mathcal A(a,t,w)-\mathcal M(a,w)
=\frac1f\sum_{g\notin\Phi}
\sum_{\chi\in\widehat\Gamma}
\overline{\chi(g)}\,w_{\rho_t(g)}
G_\chi\!\left(\frac a{w_{\rho_t(g)}}\right).
}
\label{car-balance}
\end{equation}

Esta es B, proposición 3.1 y ecuación (3.4), p. 21. La corrección principal es independiente de $t$, porque solo depende de la suma de los logaritmos de todos los pesos. La prueba algebraica admite $\ell_\kappa$ complejo. B impone $\ell_\kappa\in\R$ para que $\mathcal M$ sea real; bajo esa hipótesis, el lado derecho de \eqref{car-balance} también es real, aunque no se hayan supuesto relaciones de conjugación entre las otras constantes $\ell_\chi$.

En la aplicación, $\kappa$ corresponde a un carácter primitivo cuadrático real y $\ell_\kappa$ es real. La existencia de $L'(1,\chi^*)/L(1,\chi^*)$ para $\chi\ne1$ utiliza la no anulación clásica de los valores $L(1,\chi^*)$. Ese es el único punto de no anulación necesario para interpretar las constantes de esta identidad; no se requiere una hipótesis sobre los ceros en el interior de la franja crítica. Es la precisión de B, observación 3.2, p. 21.

\subsubsection*{Comprobación cuando todos los pesos son iguales}

Si $w_\rho=w$ para todo $\rho$, el promedio \eqref{car-promedio-clase} permite simplificar \eqref{car-suma-ponderada} a

\begin{equation}
\mathcal A(a,t,w)
=\sum_{n\ge1}\Lambda(n)
\bigl(1-a_\kappa(n)\bigr)(a/n-w)_+.
\label{car-pesos-iguales}
\end{equation}

En una potencia de primo, $a_\kappa(p^j)$ pertenece a $\{1,-1,0\}$, por lo que su coeficiente en \eqref{car-pesos-iguales} es respectivamente $0$, $2$ o $1$. Esto concuerda con \eqref{car-suma-positiva}: si $\kappa$ es trivial sobre $I_p$, la clase $jF_p+I_p$ está completamente dentro o completamente fuera de $\Phi$; si no lo es, exactamente la mitad de sus $|I_p|$ elementos queda fuera. Los factores correspondientes son $0$, $2|I_p|/|I_p|$ y $2(|I_p|/2)/|I_p|$.

En este caso, \eqref{car-balance} se reduce a

\[
\mathcal A(a,t,w)-
a\bigl(\log(a/w)-1-\gamma_E+\ell_\kappa\bigr)
=w\bigl(G_1(a/w)-G_\kappa(a/w)\bigr).
\]

Esta especialización comprueba los signos y muestra por qué la variación de los pesos tiene un papel propio en la selección posterior de otros caracteres.

\subsection{Un ejemplo que muestra el conjugado y los factores de inercia}

Tomemos

\[
\Gamma=(\Z/2\Z)\times(\Z/3\Z),
\qquad \kappa(s,r)=(-1)^s,
\qquad \Phi=\{0\}\times\Z/3\Z.
\]

Entonces $f=3$. Con $\omega=e^{2\pi i/3}$, los seis caracteres son

\[
\chi_{a,b}(s,r)=(-1)^{as}\omega^{br},
\qquad a\in\{0,1\},\quad b\in\{0,1,2\}.
\]

Para un primo fijado, usemos el representante $\widetilde F=(1,1)$ y consideremos primero el exponente $j=1$. Escribamos

\begin{equation}
K_I(g)=\frac13\sum_{a=0}^1\sum_{b=0}^2
\overline{\chi_{a,b}(g)}\,
\frac1{|I|}\sum_{h\in I}\chi_{a,b}(\widetilde F+h).
\label{car-nucleo-ejemplo}
\end{equation}

Los puntos de $\Gamma\setminus\Phi$ son $(1,0),(1,1),(1,2)$. La siguiente tabla se obtiene de la identidad \eqref{car-frobenius}:

\begin{center}
\small
\begin{tabularx}{\textwidth}{@{}lYccc@{}}
\toprule
Inercia $I$ & Caracteres que sobreviven & $K_I(1,0)$ & $K_I(1,1)$ & $K_I(1,2)$ \\
\midrule
$\{0\}$ & Los seis & $0$ & $2$ & $0$ \\
$\{0\}\times\Z/3\Z$ & Los de $b=0$ & $2/3$ & $2/3$ & $2/3$ \\
$\Z/2\Z\times\{0\}$ & Los de $a=0$ & $0$ & $1$ & $0$ \\
\bottomrule
\end{tabularx}
\end{center}

Podemos comprobarla directamente. Si $I=\{0\}$ y $g=(1,r)$, el producto de los valores de los caracteres es $\omega^{b(1-r)}$, y

\[
K_{\{0\}}(1,r)
=\frac23\sum_{b=0}^2\omega^{b(1-r)}
=2\mathbf1_{\{r=1\}}.
\]

Si $I=\{0\}\times\Z/3\Z$, la media elimina los caracteres con $b\ne0$. Los dos restantes aportan $1/3$ cada uno, lo que da $2/3$ en los tres puntos. Si $I=\Z/2\Z\times\{0\}$, solo quedan los caracteres con $a=0$ y la suma es $(1/3)\sum_b\omega^{b(1-r)}=\mathbf1_{\{r=1\}}$.

El conjugado es visible en este ejemplo. Si se omite en el primer caso y se toma $g=\widetilde F$, el resultado pasa a ser

\[
\frac23\sum_{b=0}^2\omega^{2b}=0,
\]

aunque la condición de Frobenius se cumple y el valor correcto es $2$. Este es un ejemplo algebraico propio para comprobar la normalización; no afirma que las tres inercias sean datos simultáneos de un cuerpo concreto.

También permite comparar los exponentes. En el caso no ramificado, $j\widetilde F=(j\bmod2,j\bmod3)$ queda fuera de $\Phi$ exactamente cuando $j$ es impar, y el coeficiente de la contribución local es $2\log p$. Con $I=\Z/2\Z\times\{0\}$, cada clase contiene un único punto exterior $(1,j\bmod3)$ y su coeficiente es $\log p$ para \textbf{todo entero} $j\ge1$. Esto explica el factor que A utiliza en el primo ramificado distinguido $p=d$, donde la inercia tiene orden dos y sale de $\Phi$; véase A, p. 15.

\subsection{La selección finita que prepara la combinación \texorpdfstring{$H_\psi$}{H psi}}

La identidad principal contiene todos los caracteres de $\Gamma$ y argumentos que dependen de los pesos. Para estudiar la selección posterior, aislamos ahora un cálculo finito en el grupo utilizado por A:

\[
\Gamma=(\Z/2\Z)^2\times A,
\qquad \kappa(\alpha,\beta,a)=(-1)^\alpha,
\qquad c=(1,1,0),
\]

donde $A$ es un grupo abeliano finito. Aquí $A$ designa un grupo; $\mathcal A$ sigue designando la suma aritmética. Sean $A_0\le A$, $n=|A|$ y $o=[A:A_0]$. Entonces $f=2n$. Identificamos $u\in A$ con $(0,0,u)\in\Gamma$.

Definamos

\[
\Psi=\{\chi\in\widehat\Gamma:\chi(0,0,u)=1\text{ para todo }u\in A_0\}.
\]

Estos caracteres descienden exactamente al cociente $\Gamma/A_0$, cuyo orden es $4o$. Por el conteo de caracteres ya demostrado,

\begin{equation}
|\Psi|=4o.
\label{car-caracteres-supervivientes}
\end{equation}

En la aplicación, un subconjunto $K\subset\Psi$ de tamaño $o$ representa las potencias del carácter original $\psi$. Los otros $3o$ caracteres son las tres familias obtenidas al multiplicarlas por los caracteres cuadráticos auxiliares. Esta identificación concreta corresponde a A, p. 13; depende de la construcción del cuerpo y de su corrección de paridad. El promedio por sí mismo retiene $\Psi$ completo.

\subsubsection*{El promedio sobre $A_0$ y el factor $\delta_2/2$}

Sean $T_2\subset A_0$ no vacío, $b>0$ y $t\in A$. La relación

\[
\rho_t(1,0,s+u-t)=(0,1,s+u)
\]

identifica los índices que A llama lugares de peso grande. Para cada carácter, la ortogonalidad sobre $A_0$ da

\begin{equation}
\frac1{|A_0|}\sum_{u\in A_0}\sum_{s\in T_2}
\overline{\chi(1,0,s+u-t)}
=
\begin{cases}
|T_2|\overline{\chi(1,0,-t)},&\chi\in\Psi,\\
0,&\chi\notin\Psi.
\end{cases}
\label{car-promedio-subgrupo}
\end{equation}

Para demostrarlo, se factoriza el carácter en las contribuciones de $(1,0,-t)$, $(0,0,s)$ y $(0,0,u)$. La media sobre $u$ se anula fuera de $\Psi$. Dentro de $\Psi$, tanto $s$ como $u$ pertenecen a $A_0$ y sus factores valen uno.

Si $\delta_2=|T_2|/n$, el factor que acompaña a esta media en la identidad de caracteres es

\begin{equation}
\frac{b|T_2|}{f}=\frac{b\delta_2}{2}.
\label{car-factor-medio}
\end{equation}

Esta es la procedencia del $2$ que aparece al despejar el promedio normalizado en A, proposición 4.3. El cálculo supone que lo que multiplica al carácter ya es independiente de $u$ y $s$. En la aplicación, esa independencia se obtiene después de reemplazar los pesos correspondientes por $b$ y acotar el cambio. La ortogonalidad no autoriza a extraer de una suma un factor que todavía depende de la traslación.

\subsubsection*{Por qué aparece $1/(3o+1)$}

Pongamos $Q=A/A_0$, de orden $o$, y definamos una distribución de probabilidad en sus clases:

\begin{equation}
q(C)=
\begin{cases}
4/(3o+1),&C=0,\\
3/(3o+1),&C\ne0.
\end{cases}
\label{car-distribucion}
\end{equation}

Sus valores suman uno, pues $4+3(o-1)=3o+1$. Para el carácter trivial de $Q$, la media es uno. Para cualquier carácter no trivial $\lambda\in\widehat Q$, la ortogonalidad da

\begin{equation}
\sum_{C\in Q}q(C)\lambda(C)
=\frac{3\sum_C\lambda(C)+\lambda(0)}{3o+1}
=\frac1{3o+1}.
\label{car-sesgo}
\end{equation}

Así, un pequeño exceso de masa en la clase cero produce el mismo coeficiente real positivo para todos los caracteres no triviales del cociente. A construye un conjunto $T_1\subset A$ cuya distribución entre las clases aproxima \eqref{car-distribucion}, con una condición adicional sobre órbitas. La existencia simultánea de ambas propiedades se utiliza como entrada de la construcción; el cálculo de Fourier \eqref{car-sesgo} no demuestra la condición de órbitas. Véanse A, (4.4)–(4.5), p. 14.

La estabilidad de \eqref{car-sesgo} es elemental. Si $r=|T_1|>0$, $N_C=|T_1\cap C|$ y $|N_C/r-q(C)|\le\epsilon_T$ para todas las clases, entonces

\begin{equation}
\left|\frac1r\sum_{t\in T_1}\lambda(t)-\sum_Cq(C)\lambda(C)\right|
\le\sum_C\left|\frac{N_C}{r}-q(C)\right|
\le o\epsilon_T.
\label{car-error-coeficientes}
\end{equation}

La primera desigualdad usa $|\lambda(C)|=1$. Esta cota explica cómo se controla la diferencia entre el coeficiente objetivo y el coeficiente obtenido de un conjunto finito.

\subsubsection*{Una desigualdad finita para separar los caracteres retenidos}

Supongamos que $K$ contiene al carácter trivial, que $\chi(1,0,0)=1$ para $\chi\in K$, y que la restricción a $A$, descendida a $Q$, define una biyección $K\to\widehat Q$. Estas propiedades están realizadas por las potencias de $\psi$ en el cuerpo construido por A. Fijemos números complejos $Z_\chi$, independientes de $t,u,s$, y escribamos

\begin{equation}
c_\chi=\overline{\chi(1,0,0)}\,
\frac1r\sum_{t\in T_1}\chi(0,0,t),
\qquad
\mathcal F(Z)=\sum_{\chi\in\Psi}c_\chi Z_\chi.
\label{car-funcional}
\end{equation}

La fórmula \eqref{car-promedio-subgrupo} muestra exactamente que

\begin{equation}
\mathcal F(Z)
=\frac{2}{b\delta_2 r}\sum_{t\in T_1}
\frac1{|A_0|}\sum_{u\in A_0}
\frac b f\sum_{s\in T_2}\sum_{\chi\in\widehat\Gamma}
\overline{\chi(1,0,s+u-t)}\,Z_\chi.
\label{car-promedio-completo}
\end{equation}

Se tiene $c_1=1$ y $|c_\chi|\le1$ para todo $\chi\in\Psi$. Si, además,

\[
|c_\chi-c_0|\le\rho\quad(\chi\in K\setminus\{1\}),
\qquad c_0=\frac1{3o+1},\quad\rho\ge0,
\]

entonces

\begin{equation}
\Ree Z_1+c_0\sum_{\chi\in K\setminus\{1\}}\Ree Z_\chi
\le\Ree\mathcal F(Z)
+\rho\sum_{\chi\in K\setminus\{1\}}|Z_\chi|
+\sum_{\chi\in\Psi\setminus K}|Z_\chi|.
\label{car-extraccion}
\end{equation}

En efecto, para un carácter retenido se escribe $c_\chi=c_0+e_\chi$ y se usa $\Ree(e_\chi Z_\chi)\ge-\rho|Z_\chi|$. Para los restantes, $\Ree(c_\chi Z_\chi)\ge-|Z_\chi|$. Sumando y trasladando estos errores se obtiene \eqref{car-extraccion}.

Cuando $Z_\chi=G_\chi(x)$ y $K$ se identifica con los caracteres primitivos que inducen las potencias de $\psi$, el lado izquierdo es

\[
H_\psi(x)=G_1(x)+\frac1{3o+1}
\sum_{j=1}^{o-1}\Ree G_{\psi^{[j]}}(x).
\]

Para $o=1$, la suma retenida no trivial es vacía; permanecen el carácter trivial y los tres caracteres auxiliares de $\Psi\setminus K$. La fórmula \eqref{car-extraccion} exhibe los dos errores que la selección finita deja abiertos: la aproximación de coeficientes y los caracteres auxiliares supervivientes.

\subsection{Datos aritméticos y estimaciones que aún deben incorporarse}
\label{car-entradas}

La prueba anterior tiene un núcleo exacto y una aplicación que requiere datos adicionales. Conviene fijar su relación para continuar el estudio sin atribuir a la ortogonalidad una estimación que no proporciona.

\begingroup
\raggedright
\begin{description}
\item[$I_p$, $F_p$ y carácter primitivo asociado.] Interpretar $a_\chi(n)$ como coeficientes de una función $L$ de Dirichlet. Se explicó el criterio local; la identificación aritmética clásica se utiliza como entrada.

\item[$\ell_\chi=L'(1,\chi^*)/L(1,\chi^*)$.] Recuperar las mismas sumas corregidas de la nota analítica. Se utiliza la no anulación clásica en $1$ para caracteres no principales.

\item[Cuerpo con $\Gamma=(\Z/2\Z)^2\times A$ y familia $K$.] Alojar las potencias de $\psi$, incluidas las impares, y sus twists cuadráticos. Se identificaron las condiciones usadas; la construcción completa se toma de A, pp. 12–13.

\item[Conjuntos $T_1,T_2$ y pesos.] Dar coeficientes próximos a \eqref{car-sesgo} y controlar simultáneamente las órbitas. Se demostraron sus consecuencias de Fourier; su elección aritmética y combinatoria sigue por desarrollar.

\item[Sustitución de pesos y argumentos.] Pasar de $wG_\chi(a_t/w)$ a $bG_\chi(x)$. Requiere las estimaciones de A, (4.7)–(4.9), pp. 15–16.

\item[Cota de $\mathcal A-\mathcal M$.] Acotar el promedio que aparece en \eqref{car-promedio-completo}. Depende de la comparación con periodos, normas y pendientes; continúa pendiente.

\item[Cotas para $\Psi\setminus K$.] Controlar el último término de \eqref{car-extraccion}. Son las estimaciones para los caracteres auxiliares; la ortogonalidad no los elimina.

\end{description}
\par
\endgroup

En la sustitución de argumentos hay una precisión cuantitativa útil. Si $a_t/b=x+e_t$ y $w=b+d$ con $0\le d\le h$, entonces

\begin{equation}
\frac{a_t}{w}-x=\frac{e_t-xd/b}{1+d/b},
\qquad
\left|\frac{a_t}{w}-x\right|\le |e_t|+\frac{xh}{b}.
\label{car-cambio-argumento}
\end{equation}

Por ello, conocer $h/b$ pequeño exige todavía compararlo con la escala $1/x$ para obtener un error absoluto acotado. A controla primero el cambio de $wG_\chi(a_t/w)$ al sustituir $w$ por $b$, y después usa $|a_t/b-x|\le1$ para cambiar el argumento. B, p. 22, describe las dos escalas de pesos; su expresión «próximo a $x$» debe leerse junto con ese control cuantitativo.

Finalmente, la cota geométrica que A, proposición 4.3, usa como hipótesis vale para \textbf{cada} traslación $u\in A_0$, con los mismos datos de la construcción. Para promediarla hay que conservar este cuantificador o demostrar una cota equivalente para el promedio pertinente. La identidad de caracteres determina qué exceso se debe acotar, pero no produce su cota superior.

\subsection{Concordancia con Lean y alcance de la revisión}
\begingroup
\small
\raggedright

Se consultó el espejo \href{https://github.com/CocoSgt/riemann-hypothesis-lean/tree/7cbf7a0c32fca1badf05931879a37c3f8a2b03b4}{CocoSgt/riemann-hypothesis-lean}, fijado en \texttt{7cbf7a0c32fc}, la misma revisión utilizada en la nota de ventanas. Los archivos de configuración declaran Lean \texttt{4.33.1} y Mathlib \texttt{0df444a360ea}.

La declaración central es \href{https://github.com/CocoSgt/riemann-hypothesis-lean/blob/7cbf7a0c32fca1badf05931879a37c3f8a2b03b4/src/Development/AutomorphicArithmetic/DedekindLFunctions/LandauTheorems/ChBalance.lean}{Deduction.ch\_balance}. Su conclusión es la igualdad compleja \texttt{ASK - MAIN = De}. La función \texttt{ell} puede tomar valores complejos arbitrarios; en particular, \texttt{ch\_balance} no exige que \texttt{ell chC} sea real. La interpretación como derivadas logarítmicas de caracteres primitivos se incorpora mediante \texttt{HostAddCharCompatible} y \texttt{prim\_read}.

\begingroup
\raggedright
\begin{description}
\item[Carácter trivial $1$.] $0:\texttt{AddChar}\,(\texttt{E0}\times H)\,\C$. El cero de esta notación representa el carácter constantemente igual a uno, no la función nula.

\item[Datos de inercia y Frobenius.] \texttt{ArithData}, \texttt{cosetMean}, \texttt{primVal}. El producto sobre factores primos extiende las medias de inercia a los enteros positivos.

\item[$\Phi$, $\kappa$, $\rho_t$.] \texttt{InPhi}, \texttt{PhiSet}, \texttt{chC}, \texttt{partner}. El grupo de la interfaz formal es $(\Z/2\Z)^2\times H$.

\item[Ortogonalidad y promedio de clase.] \texttt{PB.ortho}, \texttt{orthogonality\_ramified}, \texttt{colour\_sum}. Se conserva el conjugado y el factor $2/|I_p|$; en la clase exterior sobreviven $1$ y $\kappa$.

\item[Sumas corregidas y término principal.] \texttt{GPsi}, \texttt{MAIN}, \texttt{De}. Coinciden los signos, la corrección $+a\ell_\kappa$ y los argumentos $a/w$.

\item[Identidad principal.] \texttt{ch\_balance}. Se supone \texttt{HostFactsS} y un corte finito suficiente.

\item[Identificación primitiva y combinación retenida.] \texttt{prim\_read}, \texttt{gmA\_eq}. Enlazan las sumas abstractas con los caracteres primitivos y con la combinación del nivel.

\item[Coeficientes y extracción posterior.] \texttt{tops\_sum}, \texttt{ch\_core}. Se localizaron el coeficiente objetivo y la interfaz cuantitativa con sus hipótesis; el exceso geométrico sigue siendo una entrada.

\end{description}
\par
\endgroup

Los nombres de esta concordancia se abrevian por legibilidad. \texttt{InPhi} y \texttt{PhiSet} pertenecen a \texttt{Formalism.D5}; los objetos de sumas y ortogonalidad pertenecen a \texttt{Formalism.LFunctionChainParameters.characterSumData}; \texttt{ch\_balance} y \texttt{ch\_core} pertenecen a \texttt{Deduction}; \texttt{prim\_read}, \texttt{gmA\_eq} y \texttt{tops\_sum} pertenecen a \texttt{UpperTheory.Comparison}$_2$. El manifiesto conserva los nombres completos.

B enuncia la identidad finita para cualquier grupo abeliano finito con un carácter cuadrático no trivial. La entrada \texttt{ch\_balance} consultada utiliza el grupo $(\Z/2\Z)^2\times H$ y \texttt{HostFactsS}. Estas hipótesis identifican \texttt{ASK} con su forma concreta: contribuciones de primos inertes en el subcampo cuadrático distinguido, contribuciones nulas cuando la inercia y Frobenius quedan en $\Phi$, y la contribución del primo ramificado distinguido. La demostración abstracta de esta nota no requiere construir ese cuerpo; el cotejo de Lean corresponde a la interfaz específica descrita.

\texttt{prim\_read} identifica \texttt{GPsi} con la suma corregida del carácter primitivo bajo los datos del host y \texttt{HostAddCharCompatible}. \texttt{gmA\_eq}, bajo \texttt{EAtH} y esa compatibilidad, identifica la suma retenida con \texttt{Level.Gm}. Para \texttt{Level.cyc}\,$\psi$, el campo \texttt{psK} es \texttt{none} y se obtiene exactamente $H_\psi$, con peso $1$ para el principal y $1/(3o+1)$ para las otras primitivas retenidas. Los niveles generales permiten un carácter excepcional de peso $1$. Estas declaraciones establecen identidades entre las sumas; la cota unilateral y la cota del exceso geométrico requieren las entradas adicionales señaladas en §\ref{car-entradas}.

En la interfaz de extracción leída, los pesos tienen la forma $w(u,j,\rho)$, independiente de $t$, y el ancho de la banda es $\eta+1$, con $\eta\ge0$. A, proposición 4.3, permite además dependencia en $t$ y escribe el ancho como $h$; la especialización formal corresponde a $h=\eta+1$. El cotejo identifica los pasos finitos y su ensamble, sin presentar la proposición de A como transcripción literal del enunciado final \texttt{ch\_core}.

Además, Lean lleva un entero de corte $N$ y exige $a/w_\rho\le N$ para $\rho\in\Phi$. Nuestra notación de sumas sobre todos los enteros es equivalente porque los sumandos se anulan fuera de ese rango. Un entero $N\ge a/w_{\min}$ sirve para todas las posiciones. También acota los primos y exponentes relevantes: si $p^j$ contribuye, entonces $p\le p^j<N$ y $j\le p^j<N$.

El registro acumulativo del repositorio conserva las rutas, declaraciones, hashes y el alcance de lectura de los archivos recuperados. Este cotejo no incluyó compilación ni ejecución de Lean, una auditoría de todas las dependencias ni una comparación integral del espejo con la distribución original de Alpöge.

\paragraph{Estado del tramo.} Quedan demostradas la ortogonalidad finita, la identidad con inercia, la cancelación del término principal, el ejemplo local y las consecuencias finitas de los dos promedios. La estimación aritmética completa y la cota geométrica permanecen identificadas como entradas. El siguiente desarrollo será la desigualdad de marcos y su paso a pendientes medias, que participa en el control de ese exceso. Los pendientes se mantienen en el registro de trabajo del cuaderno.
\par
\endgroup

\clearpage
\section{Marcos aproximados y pendientes medias}
\label{sec:marcos}

\fuente{A, §5, Lema~5.1, Proposición~5.2 y Teorema~5.3,
pp.~17--20; B, §§4.3--4.4, Lema~4.2 y Proposición~4.3,
pp.~39--44. Reconstruimos la desigualdad lineal y su integración.
La existencia de la filtración de Harder--Narasimhan se utiliza como
resultado clásico citado.}

Las aplicaciones integrales de la Nota~\ref{sec:proyecciones} parten de
retículos pequeños y llegan a un mismo retículo receptor. Para comparar
sus pendientes medias interesa saber cuántas direcciones del receptor
siguen cubiertas cuando se retienen solo ciertas partes de las fuentes.
Una familia muy numerosa puede concentrarse en una sola dirección.
La información pertinente será una estimación colectiva de trazas.

El argumento se desarrolla en dos niveles. Primero se estima la
codimensión del espacio generado por imágenes de subespacios. Después
se eligen esos subespacios mediante umbrales de pendiente, se aplica la
estimación para cada umbral y se integra el déficit resultante. Las
condiciones que hacen posible aplicar este mecanismo a los periodos
automorfos quedan identificadas al final.

\subsection{Operadores de Gram y distribución de las imágenes}

Sean $V_i$ y $H$ espacios de Hilbert complejos de dimensión finita, y
$B_i:V_i\to H$ una familia finita de aplicaciones lineales. Escribimos
\[
 r_i=\dim V_i,\qquad R=\sum_i r_i>0,\qquad D=\dim H>0.
\]
El adjunto $B_i^*$ se caracteriza por la identidad del producto hermitiano.
Para un operador $A$, su norma de Hilbert--Schmidt satisface
\[
 \|A\|_{\mathrm{HS}}^2=\operatorname{Tr}(A^*A).
\]
En bases ortonormales es la suma de los cuadrados de los valores absolutos
de las entradas. La desigualdad de Cauchy--Schwarz para este producto
escalar da
$|\operatorname{Tr}(A^*C)|\le\|A\|_{\mathrm{HS}}\|C\|_{\mathrm{HS}}$.
Un proyector ortogonal de rango $k$ tiene norma de Hilbert--Schmidt
$\sqrt{k}$.

Elegimos números $c_i>0$ y definimos
\begin{equation}
 A_i=c_iB_i^*B_i,\qquad \Delta_i=A_i-I_{V_i},\qquad
 Q=\sum_i c_iB_iB_i^*.
 \label{eq:marcos-operadores}
\end{equation}
Los operadores $A_i$ y $Q$ son autoadjuntos semidefinidos positivos;
$\Delta_i$ puede tener autovalores de ambos signos. La identidad
de trazas para aplicaciones entre espacios de dimensiones distintas da
\[
 T:=\operatorname{Tr}Q=\sum_i\operatorname{Tr}A_i.
\]

Las dos hipótesis del marco son
\begin{align}
 \sum_i\|\Delta_i\|_{\mathrm{HS}}^2&\le eR,
       &&0\le e\le\tfrac14,
       \label{eq:marcos-error-fuentes}\\
 \operatorname{Tr}(Q^2)&\le(1+\eta)\frac{T^2}{D},
       &&\eta\ge0.
       \label{eq:marcos-error-destino}
\end{align}
La primera compara las métricas de las fuentes con las inducidas por
los mapas ponderados. La segunda limita la dispersión de los autovalores
del operador colectivo. En efecto,
\begin{equation}
 \left\|Q-\frac TD I_H\right\|_{\mathrm{HS}}^2
 =\operatorname{Tr}(Q^2)-\frac{T^2}{D}
 \le\eta\frac{T^2}{D}.
 \label{eq:marcos-varianza}
\end{equation}
Si $e=\eta=0$, cada $A_i$ es la identidad y $Q=(R/D)I_H$.
La escala escalar del marco es $R/D$; puede ser distinta de $1$.

\subsection{La desigualdad de codimensión}

\begin{proposicion}[Marco aproximado]
Supongamos \eqref{eq:marcos-error-fuentes} y
\eqref{eq:marcos-error-destino}. Si $N_i\subseteq V_i$ son
subespacios y $\varepsilon\ge0$ satisface
$\sum_i\operatorname{codim}_{V_i}N_i\le\varepsilon R$, entonces
\begin{equation}
 \operatorname{codim}_H\left(\sum_i B_iN_i\right)
 \le(\eta+4\varepsilon+4\sqrt e)D.
 \label{eq:marcos-codimension}
\end{equation}
\label{prop:marco-lineal}
\end{proposicion}

\begin{proof}
Como $T-R=\sum_i\operatorname{Tr}\Delta_i$, dos aplicaciones de
Cauchy--Schwarz dan
\begin{align*}
 |T-R|&\le\sum_i\sqrt{r_i}\|\Delta_i\|_{\mathrm{HS}}\\
 &\le\sqrt R\left(\sum_i\|\Delta_i\|_{\mathrm{HS}}^2\right)^{1/2}
 \le R\sqrt e.
\end{align*}
Por tanto,
\begin{equation}
 T\ge(1-\sqrt e)R\ge R/2>0.
 \label{eq:marcos-traza-positiva}
\end{equation}

Sea $S=\sum_iB_iN_i$. Denotemos por $P$ el proyector ortogonal
sobre $S^\perp$ y por $P_i$ el proyector sobre $N_i^\perp$.
La aplicación $PB_i$ se anula en $N_i$, de modo que $PB_i=PB_iP_i$.
Una proyección ortogonal no aumenta las normas de los vectores; aplicando
ese hecho a una base ortonormal de $V_i$ obtenemos
\begin{align}
 \operatorname{Tr}(PQ)
 &=\sum_i c_i\|PB_i\|_{\mathrm{HS}}^2\notag\\
 &\le\sum_i c_i\|B_iP_i\|_{\mathrm{HS}}^2
 =\sum_i\operatorname{Tr}(P_iA_i)\notag\\
 &=\sum_i\operatorname{rk}P_i+
       \sum_i\operatorname{Tr}(P_i\Delta_i).
 \label{eq:marcos-traza-omitida}
\end{align}
Las trazas del primer miembro y de $P_iA_i$ son reales y no negativas:
por ciclicidad, coinciden con las de $PQP$ y $P_iA_iP_i$.
Esto no exige que los productos $PQ$ o $P_iA_i$ sean autoadjuntos.

Para el término de error,
\begin{align*}
 \left|\sum_i\operatorname{Tr}(P_i\Delta_i)\right|
 &\le\sum_i\sqrt{\operatorname{rk}P_i}\|\Delta_i\|_{\mathrm{HS}}\\
 &\le\left(\sum_i\operatorname{rk}P_i\right)^{1/2}
       \left(\sum_i\|\Delta_i\|_{\mathrm{HS}}^2\right)^{1/2}
 \le R\sqrt e,
\end{align*}
pues $\sum_i\operatorname{rk}P_i\le R$. La hipótesis de codimensión
controla el término de identidad de \eqref{eq:marcos-traza-omitida}, y da
\begin{equation}
 0\le\operatorname{Tr}(PQ)\le(\varepsilon+\sqrt e)R.
 \label{eq:marcos-control-omitido}
\end{equation}

Ponemos $z=\operatorname{Tr}(PQ)/T$. La positividad de las trazas
sobre $S$ y $S^\perp$ implica $0\le z\le1$;
\eqref{eq:marcos-traza-positiva} y \eqref{eq:marcos-control-omitido}
implican $z\le2(\varepsilon+\sqrt e)$. Aplicando Cauchy--Schwarz
a $I_H-P$ y $Q$,
\[
 (1-z)^2T^2
 =\bigl(\operatorname{Tr}((I_H-P)Q)\bigr)^2
 \le\dim S\,\operatorname{Tr}(Q^2)
 \le\dim S\,(1+\eta)T^2/D.
\]
Podemos cancelar $T^2>0$. Además,
\[
 \frac{(1-z)^2}{1+\eta}\ge1-\eta-2z,
\]
porque al multiplicar por $1+\eta$ la diferencia de ambos miembros
es $z^2+\eta^2+2\eta z\ge0$. Por consiguiente,
\[
 \frac{\dim S}{D}\ge1-\eta-4\varepsilon-4\sqrt e,
\]
que equivale a \eqref{eq:marcos-codimension}.
\end{proof}

Esta es la prueba de A, Proposición~5.2, desarrollada con los operadores
de compresión explícitos. B, Lema~4.2, conserva primero la cota más fina
$R(\varepsilon+\sqrt{\varepsilon e})$ y compara la traza omitida
con \eqref{eq:marcos-varianza}; llega a la misma constante.
Para $\varepsilon\ge1$ la conclusión es inmediata de
$\operatorname{codim}S\le D$. Los mapas pueden tener núcleo y algunos
$N_i$ pueden ser cero. Cuando el coeficiente de la derecha es mayor
o igual que $1$, la desigualdad sigue siendo válida, aunque no mejora
esa cota elemental.

\subsection{Ejemplos de concentración y redundancia}

Tomemos $V_i=\C$, $H=\C^D$, $c_i=1$ y $B_i(z)=ze_1$ para
$1\le i\le n$, donde $e_1$ pertenece a la base ortonormal habitual.
Entonces $A_i=1$, así que $e=0$ es admisible, pero
\[
 Q=nP_{\C e_1},\qquad T=n,\qquad
 \operatorname{Tr}(Q^2)=n^2.
\]
La segunda hipótesis exige $\eta\ge D-1$. Aumentar $n$ conserva
esa exigencia: todas las imágenes siguen en la misma recta.

En cambio, si se toman dos copias de cada aplicación
$\C\to\C^D$, $z\mapsto ze_j$, entonces $R=2D$ y $Q=2I_H$.
Se tiene $e=\eta=0$. Hay redundancia: la aplicación agregada
$\bigoplus_iV_i\to H$, $(v_i)_i\mapsto\sum_iB_iv_i$, tiene
núcleo de dimensión $D$. Ese núcleo no es el núcleo de un mapa
individual y es compatible con un marco exacto.

Con una sola copia de cada coordenada y $D=10$, suprimir una fuente
de dimensión uno da $\varepsilon=1/10$. La imagen retenida tiene
codimensión $1$, mientras \eqref{eq:marcos-codimension} garantiza
codimensión a lo sumo $4$. El ejemplo muestra el carácter uniforme
de la constante, que no pretende ser óptima en cada familia.

\subsection{Normalización aritmética de las pendientes}

Sea $K$ un cuerpo de números, de grado $d_K$, y sea $E$ un
retículo hermitiano: un $\OK$-módulo proyectivo de rango finito positivo con
métricas compatibles en sus realizaciones arquimedianas. Para una
sección racional no nula $v$ de $\det E$, el grado normalizado es
\begin{equation}
 \widehat{\deg}_{\mathrm{Ar}}(E)=\frac1{d_K}
 \left(\sum_{\mathfrak p}\operatorname{ord}_{\mathfrak p}(v;\det E)
                 \log N\mathfrak p
       -\sum_{\sigma:K\hookrightarrow\C}\log\|v\|_\sigma\right).
 \label{eq:marcos-grado}
\end{equation}
La primera suma recorre los primos finitos. El orden se mide respecto
de un generador local y puede ser negativo. La fórmula del producto
hace que el resultado sea independiente de $v$; no se necesita una
base integral global de $E$.

Usamos la normalización de A y B:
\begin{equation}
 \mu(E)=\frac{2\widehat{\deg}_{\mathrm{Ar}}(E)}{\operatorname{rk}E},
 \qquad
 h(B)=\frac2{d_K}\sum_{\sigma:K\hookrightarrow\C}
                         \log\|B_\sigma\|_{\mathrm{op}}\quad(B\ne0).
 \label{eq:marcos-normalizacion}
\end{equation}
La suma comprende todas las inmersiones, contando las dos de cada par
conjugado. Para $K$ totalmente complejo es equivalente sumar
$\log\|B_v\|_{\mathrm{op}}^2$ una vez por lugar complejo y multiplicar
por $2/d_K$, como en B. El número $h(B)$ puede ser negativo.
Sobre $\Q$, $\mu(E)=-2\log\operatorname{covol}(E)/\operatorname{rk}E$,
el doble de la pendiente introductoria.

Si $K$ es totalmente complejo, $r_2=d_K/2$ y $r=\operatorname{rk}E$,
la relación con el covolumen euclídeo de Minkowski es
\begin{equation}
 \mathcal H(E)=2^{r_2r}|\Delta_K|^{-r/2}
                       \operatorname{covol}_{\R}(E),
 \qquad \mu(E)=-\frac2{d_Kr}\log\mathcal H(E).
 \label{eq:marcos-covolumen}
\end{equation}
El factor finito normaliza el covolumen de $\OK$; con estas
convenciones $\mu(\OK^r)=0$ para las métricas habituales.
Las fórmulas \eqref{eq:marcos-grado}--\eqref{eq:marcos-covolumen}
son las de A, (5.1)--(5.3), pp.~17--18, y B, p.~39.

Un subretículo $M\subset E$ es \emph{saturado} cuando $E/M$ no
tiene torsión. La saturación reemplaza $M$ por
$E\cap(M\otimes_{\OK}K)$. Si $E'\subset E$ tiene el mismo rango
$r>0$, las mismas métricas y un cociente finito de cardinal $J$, entonces
\begin{equation}
 \widehat{\deg}_{\mathrm{Ar}}(E')=
 \widehat{\deg}_{\mathrm{Ar}}(E)-\frac{\log J}{d_K},
 \qquad \mu(E')=\mu(E)-\frac{2\log J}{d_Kr}.
 \label{eq:marcos-coste-indice}
\end{equation}
Ampliar el retículo invierte el signo de esa diferencia. Esta es la
contribución que conecta los denominadores de la
Nota~\ref{sec:proyecciones} con las pendientes.

\subsection{La filtración de Harder--Narasimhan}

La pendiente máxima $\mu_{\max}(E)$ es el máximo de las pendientes
de los subretículos saturados no nulos. La mínima $\mu_{\min}(E)$
es el mínimo de las pendientes de los cocientes no nulos sin torsión,
provistos de las métricas cociente. Un retículo es \emph{semiestable}
cuando $\mu_{\max}(E)=\mu(E)$.

La teoría clásica proporciona una filtración canónica por subretículos
saturados
\[
 0=E_0\subset E_1\subset\cdots\subset E_\ell=E,
\]
cuyos cocientes sucesivos, con las métricas inducidas, son semiestables
y tienen pendientes estrictamente decrecientes. Repetir cada pendiente
según el rango de su cociente produce una lista
\begin{equation}
 \lambda_1(E)\ge\cdots\ge\lambda_r(E),\qquad
 \sum_{j=1}^r\lambda_j(E)=r\mu(E).
 \label{eq:marcos-lista-hn}
\end{equation}
La primera y la última son $\mu_{\max}(E)$ y $\mu_{\min}(E)$.
La existencia y las propiedades de esta filtración se reciben de la
teoría de reducción aritmética; véanse Grayson \cite{grayson-hn},
la discusión de A, p.~18, y B, p.~42. La aditividad del grado con
métricas de subespacio y cociente está recogida también en
\cite[§4.1.1, Lema~4.2]{bost-pendientes}.

Escribimos $E^{\ge t}$ para la etapa de la filtración que retiene
los cocientes de pendiente al menos $t$. Entonces
\begin{align}
 \operatorname{rk}E^{\ge t}&=\#\{j:\lambda_j(E)\ge t\},
                  \label{eq:marcos-corte-hn}\\
 \mu_{\min}(E^{\ge t})&\ge t,\qquad
 \mu_{\max}(E/E^{\ge t})<t,
                  \label{eq:marcos-extremos-hn}
\end{align}
donde cada pendiente se considera solo si el retículo correspondiente
es no nulo. La desigualdad estricta del cociente resulta de haber
retenido también los cocientes de pendiente igual a $t$.
No se elige una descomposición de $E$ en los cocientes sucesivos:
la filtración y sus rangos bastan para el argumento.

\subsection{Una aplicación integral desplaza los umbrales}

\begin{lema}[Comparación de pendientes]
Si $B:E\to F$ es una aplicación integral no nula entre retículos
hermitianos, entonces
\begin{equation}
 \mu_{\min}(E)\le\mu_{\max}(F)+h(B).
 \label{eq:marcos-morfismo-pendientes}
\end{equation}
Para todo $t\in\R$ se tiene además
\begin{equation}
 B(E^{\ge t})\subseteq F^{\ge t-h(B)}.
 \label{eq:marcos-inclusion-filtrada}
\end{equation}
\label{lem:marcos-altura}
\end{lema}

\begin{proof}
Sea $I=B(E)$ con la métrica inducida por $F$, y sea $q=\operatorname{rk}I$.
El núcleo es saturado porque $E/\ker B\simeq I$ no tiene torsión.
La aplicación inducida $\overline B:E/\ker B\to I$ es un isomorfismo
de módulos. Con la métrica cociente en el dominio, su determinante
tiene norma a lo sumo $\|B_\sigma\|_{\mathrm{op}}^q$ en cada
inmersión. La fórmula del grado, aplicada a ese determinante, da
\[
 \mu(E/\ker B)\le\mu(I)+h(B).
\]
La saturación $I^{\mathrm{sat}}$ en $F$ tiene el mismo rango y grado
mayor o igual. Por las definiciones de pendiente mínima y máxima,
\[
 \mu_{\min}(E)\le\mu(E/\ker B)
 \le\mu(I)+h(B)\le\mu(I^{\mathrm{sat}})+h(B)
 \le\mu_{\max}(F)+h(B).
\]

Para probar la inclusión, compongamos la restricción de $B$ a
$E^{\ge t}$ con la aplicación cociente
$F\to F/F^{\ge t-h(B)}$. Si la composición fuera no nula, su
altura sería a lo sumo $h(B)$, porque ni restringir el dominio con
su métrica inducida ni proyectar con la métrica cociente aumenta
la norma de operador. Las desigualdades
\eqref{eq:marcos-morfismo-pendientes} y
\eqref{eq:marcos-extremos-hn} darían
\[
 t\le\mu_{\min}(E^{\ge t})
 \le\mu_{\max}(F/F^{\ge t-h(B)})+h(B)<t,
\]
una contradicción. La composición es cero, lo que prueba
\eqref{eq:marcos-inclusion-filtrada}.
\end{proof}

Este es el Lema~5.1 de A. La comparación mediante determinantes es
una forma del método de pendientes de Bost
\cite[§4.1.4, pp.~199--200]{bost-pendientes}. Se han mantenido
las métricas cociente y la saturación para que la contribución integral
y la norma arquimediana correspondan a la misma aplicación.

La integridad es una hipótesis efectiva. En rango uno sobre $\Q$, con
retículos y métricas habituales, multiplicar por $1/p$ tiene altura
arquimediana $-2\log p$, pero ambas pendientes son cero. Aplicar
\eqref{eq:marcos-morfismo-pendientes} a ese mapa racional daría
$0\le-2\log p$, que es falso. El denominador introduce un coste
finito que la hipótesis integral evita omitir.

\subsection{La proporción de pendientes bajas en las fuentes}

Sean $E_i$ retículos de rangos $r_i>0$, $R=\sum_i r_i$ y
$h_i\in\R$. La lista de pendientes desplazadas es
\[
 a_{i,j}=\lambda_j(E_i)-h_i\qquad(1\le j\le r_i).
\]
Cada término recibe peso $1/R$. La media de esta lista es
\begin{equation}
 \frac1R\sum_{i,j}a_{i,j}
 =\sum_i\frac{r_i}{R}\bigl(\mu(E_i)-h_i\bigr).
 \label{eq:marcos-media-fuentes}
\end{equation}
Los pesos $r_i/R$ proceden de contar dimensiones. Sustituirlos por
probabilidades arbitrarias modificaría la relación con la hipótesis
de marco, que está normalizada por $R$.

\begin{lema}[Cola de una lista de pendientes]
Supongamos que la media de los $a_{i,j}$ es al menos $s$ y que
todos satisfacen $a_{i,j}\le s+u$, con $u\ge0$. Para $\tau>0$, sea
\[
 q(\tau)=\frac1R\#\{(i,j):a_{i,j}<s-\tau\}.
\]
Entonces
\begin{equation}
 q(\tau)\le\frac{u}{u+\tau}.
 \label{eq:marcos-cola-fuentes}
\end{equation}
Si $u=0$, todos los términos son $s$ y $q(\tau)=0$.
\label{lem:marcos-cola}
\end{lema}

\begin{proof}
La parte de la lista contada por $q(\tau)$ está por debajo de
$s-\tau$, y el resto está a lo sumo en $s+u$. Por tanto,
\[
 s\le\frac1R\sum_{i,j}a_{i,j}
 \le q(\tau)(s-\tau)+(1-q(\tau))(s+u)
 =s+u-q(\tau)(u+\tau).
\]
Como $u+\tau>0$, esta desigualdad implica
\eqref{eq:marcos-cola-fuentes}. Si $u=0$, una lista cuyos términos
son a lo sumo $s$ y cuya media es al menos $s$ tiene todos sus
términos iguales a $s$.
\end{proof}

La cota superior individual controla cuánto pueden compensar las
pendientes altas a las bajas. La media por sí sola no limita esa
proporción: una sola pendiente arbitrariamente grande podría compensar
muchas pendientes pequeñas.

\subsection{El teorema de transferencia de la pendiente media}

\begin{teorema}[Transferencia mediante un marco aproximado]
Sea $K$ un cuerpo de números totalmente complejo. Sean $E_i$ una
familia finita no vacía de retículos hermitianos de rangos $r_i>0$,
y $F$ otro de rango $r>0$. Sean $b_i:E_i\to F$ aplicaciones
integrales y $h_i\in\R$ tales que $h(b_i)\le h_i$ cuando $b_i\ne0$.
Escribamos $R=\sum_i r_i$. Supongamos que, para ciertos
$s\in\R$ y $u,H_0\ge0$,
\begin{align}
 \sum_i\frac{r_i}{R}\bigl(\mu(E_i)-h_i\bigr)&\ge s,
                   \label{eq:marcos-hip-media}\\
 \mu_{\max}(E_i)-h_i&\le s+u\quad\text{para todo }i,
                   \label{eq:marcos-hip-max}\\
 \mu_{\min}(F)&\ge s-H_0.
                   \label{eq:marcos-hip-min}
\end{align}
En una inmersión compleja fijada, elijamos estructuras de Hilbert
auxiliares, un espacio $\widetilde H$ de dimensión $D>0$, una
aplicación lineal sobreyectiva $\pi:\widetilde H\to F_\sigma$ y
levantamientos $\widetilde B_i:(E_i)_\sigma\to\widetilde H$ con
$\pi\widetilde B_i=(b_i)_\sigma$.
Supongamos que los levantamientos satisfacen
\eqref{eq:marcos-error-fuentes}--\eqref{eq:marcos-error-destino}
con coeficientes positivos, $0\le e\le1/4$ y $\eta\ge0$, y que
\begin{equation}
 r\ge(1-\delta)D,\qquad 0\le\delta<1.
 \label{eq:marcos-hip-proyeccion}
\end{equation}
Entonces
\begin{equation}
 \boxed{\quad
 \mu(F)\ge s-
 \frac{4u\log(1+H_0/u)+(\eta+4\sqrt e)H_0}{1-\delta}.
 \quad}
 \label{eq:marcos-transferencia}
\end{equation}
Cuando $u=0$, el término $u\log(1+H_0/u)$ se interpreta como $0$.
\label{teo:marcos-transferencia}
\end{teorema}

\begin{proof}
Fijemos $\tau>0$. En cada fuente retenemos la realización compleja
del corte de Harder--Narasimhan
\[
 N_i(\tau)=(E_i^{\ge s-\tau+h_i})_\sigma.
\]
Su codimensión cuenta las pendientes desplazadas menores que
$s-\tau$. Las hipótesis \eqref{eq:marcos-hip-media} y
\eqref{eq:marcos-hip-max}, junto con el Lema~\ref{lem:marcos-cola},
dan
\begin{equation}
 \frac1R\sum_i\operatorname{codim}N_i(\tau)
 =q(\tau)\le\frac{u}{u+\tau}.
 \label{eq:marcos-descarte-fuentes}
\end{equation}

Si $b_i\ne0$, el Lema~\ref{lem:marcos-altura} implica
\[
 b_i(E_i^{\ge s-\tau+h_i})
 \subseteq F^{\ge s-\tau+h_i-h(b_i)}
 \subseteq F^{\ge s-\tau}.
\]
La segunda inclusión usa $h_i\ge h(b_i)$ y que los cortes disminuyen
al aumentar el umbral. Para $b_i=0$ la inclusión es automática.

Aplicamos la Proposición~\ref{prop:marco-lineal} a los levantamientos
y a $N_i(\tau)$. Si $\widetilde S=\sum_i\widetilde B_iN_i(\tau)$,
entonces
\[
 \operatorname{codim}_{\widetilde H}\widetilde S
 \le\left(\eta+4\sqrt e+\frac{4u}{u+\tau}\right)D.
\]
La aplicación sobreyectiva $\pi$ induce una sobreyección
$\widetilde H/\widetilde S\to F_\sigma/\pi\widetilde S$.
Por ello proyectar no aumenta la codimensión. Además,
$\pi\widetilde S\subseteq(F^{\ge s-\tau})_\sigma$.
Al dividir por $r\ge(1-\delta)D>0$ obtenemos
\begin{equation}
 q_F(\tau):=\frac{\#\{j:\lambda_j(F)<s-\tau\}}r
 \le\frac1{1-\delta}
       \left(\eta+4\sqrt e+\frac{4u}{u+\tau}\right).
 \label{eq:marcos-cola-destino}
\end{equation}
Esta parte del razonamiento utiliza las métricas auxiliares solo para
aplicar el resultado lineal. Las alturas $h(b_i)$ se siguen calculando
con las métricas aritméticas originales y con todas las inmersiones.
La sobreyección $\pi$ no necesita ser ortogonal ni tener una norma acotada.

Por \eqref{eq:marcos-hip-min}, cada pendiente del destino satisface
$\lambda_j(F)\ge s-H_0$. Con $a_+=\max(a,0)$,
\begin{equation}
 (s-\lambda_j(F))_+
 =\int_0^{H_0}\mathbf1_{\{\lambda_j(F)<s-\tau\}}\dd\tau.
 \label{eq:marcos-capas}
\end{equation}
La igualdad se comprueba midiendo el intervalo
$0<\tau<s-\lambda_j(F)$ cuando el extremo es positivo.
Las posibles igualdades en los extremos no afectan a la integral.
Al promediar una suma finita,
\begin{equation}
 s-\mu(F)\le\frac1r\sum_j(s-\lambda_j(F))_+
 =\int_0^{H_0}q_F(\tau)\dd\tau.
 \label{eq:marcos-deficit-medio}
\end{equation}
Si $u>0$, integrar \eqref{eq:marcos-cola-destino} y usar
\[
 \int_0^{H_0}\frac{u}{u+\tau}\dd\tau
 =u\log(1+H_0/u)
\]
prueba \eqref{eq:marcos-transferencia}.
Si $u=0$, $q(\tau)=0$ para todo $\tau>0$ y la misma integración
solo conserva el término $(\eta+4\sqrt e)H_0/(1-\delta)$.
Si $H_0=0$, todas las pendientes de $F$ son al menos $s$ y la
conclusión se sigue directamente.
\end{proof}

La fórmula coincide con A, Teorema~5.3, y B, Proposición~4.3.
El número de mapas no aparece en la pérdida. La información sobre
toda la familia se reúne mediante los rangos y los dos defectos del
marco; no se introduce un término $\log|I|$ al integrar.

\subsection{Cómo la integral recupera un déficit}

Consideremos la lista de pendientes $(-2,0,1,3)$, con $s=1$ y
$H_0=3$. La proporción por debajo de $1-\tau$ y su contribución
a la integral son
\begin{center}
\begin{tabular}{@{}ccc@{}}
\toprule
Intervalo & $q_F(\tau)$ & Área \\
\midrule
$0<\tau<1$ & $1/2$ & $1/2$ \\
$1\le\tau<3$ & $1/4$ & $1/2$ \\
$\tau\ge3$ & $0$ & $0$ \\
\bottomrule
\end{tabular}
\end{center}
El área total es $1$, igual al déficit positivo medio
$(3+1+0+0)/4$. La media de la lista es $1/2$, de modo que
$s-\mu=1/2$. La diferencia entre ambas cantidades es el excedente
positivo medio $(0+0+0+2)/4=1/2$. Por eso
\eqref{eq:marcos-deficit-medio} comienza con una desigualdad.
El ejemplo ilustra la identidad de integración; por sí solo no
especifica mapas que satisfagan las hipótesis del teorema.

\subsection{El coste logarítmico y el tamaño de los errores}

\begin{lema}
Si $u,U,H_0,b,X\ge0$, $u\le U$ y $H_0\le Xb$, entonces
\begin{equation}
 u\log(1+H_0/u)\le(U+b)\log(1+X),
 \label{eq:marcos-absorcion-log}
\end{equation}
con la misma convención cuando $u=0$.
\end{lema}

\begin{proof}
Para $u>0$, la recta tangente a $\log y$ en $y=1+X$ da
\[
 \log(1+H_0/u)
 \le\log(1+X)+\frac{H_0/u-X}{1+X}.
\]
Multiplicar por $u$, usar $H_0\le Xb$ y descartar el término
no positivo $-uX/(1+X)$ produce
\[
 u\log(1+H_0/u)
 \le u\log(1+X)+\frac{Xb}{1+X}.
\]
La desigualdad $X/(1+X)\le\log(1+X)$ se deduce de que
la diferencia vale cero en $X=0$ y tiene derivada
$X/(1+X)^2\ge0$. Con $u\le U$ obtenemos
\eqref{eq:marcos-absorcion-log}. Para $u=0$ el miembro izquierdo
es cero y el derecho es no negativo.
\end{proof}

Este es el Lema~5.4 de A, p.~20. Si $X$ es una potencia fija de
$x$, su contribución tiene un factor $\log x$. Sustituirla
prematuramente por un producto $uH_0$ perdería esa información.
Los errores del marco aparecen multiplicados por $H_0$, mientras
que el descarte por pendientes se integra mediante la función
$u/(u+\tau)$. La pérdida $1/(1-\delta)$ procede de comparar
la dimensión auxiliar con el rango del receptor.

\subsection{Qué debe aportar la construcción de periodos}

La Nota~\ref{sec:proyecciones} construye formalmente aplicaciones
integrales $b_\alpha:N_\alpha\to C$ a partir de la divisibilidad
de los emparejamientos. El Teorema~\ref{teo:marcos-transferencia}
indica qué información adicional debe acreditarse para utilizarlas:

\begin{center}
\small
\begin{tabularx}{\textwidth}{@{}>{\raggedright\arraybackslash}p{0.32\textwidth}Y@{}}
\toprule
Dato de entrada & Función en el argumento \\
\midrule
Integridad de $b_\alpha$ y altura $h_\alpha$ &
Transportar el corte de pendiente sin omitir índices ni denominadores. \\
Media de $\mu(N_\alpha)-h_\alpha$ y cota de $\mu_{\max}(N_\alpha)-h_\alpha$ &
Fijar $s,u$ y controlar la fracción descartada en cada umbral. \\
Mínimo de las pendientes de $C$ &
Fijar el extremo $H_0$ de la integral del déficit. \\
Defectos $e,\eta$ del marco &
Controlar la pérdida de direcciones tras restringir simultáneamente las fuentes. \\
Rango retenido y dimensión auxiliar &
Fijar $\delta$ y comparar codimensión con proporción de pendientes del receptor. \\
\bottomrule
\end{tabularx}
\end{center}

Estas cotas tienen que referirse a los mismos mapas y a normalizaciones
compatibles. Las normas individuales no proporcionan por sí solas
la estimación colectiva de $\operatorname{Tr}(Q^2)$, como muestra
el ejemplo de concentración. Un descarte previo de fuentes que cambie
$e$ o $\eta$ tampoco se identifica con el parámetro $\delta$, que
mide una pérdida de dimensión en el destino. B, pp.~44--46, distingue
esos pasos en la aplicación a periodos.

La desigualdad abstracta queda demostrada con las entradas indicadas.
Producirlas para la familia automorfa requiere estudiar las torres theta,
las estimaciones de evaluación, las estructuras integrales y los costes
de las proyecciones de A, §§6--8, y B, §§4.1--4.5 y~5.
Esta nota no da todavía la cota del exceso geométrico de la
Nota~\ref{sec:caracteres} ni la ventana aritmética de la
Nota~\ref{sec:ventanas}.

\subsection{Concordancia localizada con Lean}
\label{subsec:marcos-lean}

El cotejo usa el espejo citado en \cite{lean}, fijado en la revisión
\href{https://github.com/CocoSgt/riemann-hypothesis-lean/tree/7cbf7a0c32fca1badf05931879a37c3f8a2b03b4}{\texttt{7cbf7a0}},
cuyo identificador completo figura en el registro.
Se leyeron catorce archivos de los módulos de Harder--Narasimhan y
de corangos. El registro conserva sus rutas, declaraciones, hipótesis,
URL fijadas y hashes. Los nombres siguientes llevan el prefijo
\texttt{Formalism.HN}, salvo que se indique otro.

\begin{enumerate}
\item \textbf{La constante del marco lineal.}
En
\href{https://github.com/CocoSgt/riemann-hypothesis-lean/blob/7cbf7a0c32fca1badf05931879a37c3f8a2b03b4/src/Formalism/HarderNarasimhan/StatementFrameF.lean}{\texttt{StatementFrameF.lean}},
\texttt{FrameLA1} define el enunciado
\eqref{eq:marcos-codimension}. Su prueba es
\href{https://github.com/CocoSgt/riemann-hypothesis-lean/blob/7cbf7a0c32fca1badf05931879a37c3f8a2b03b4/src/Formalism/HarderNarasimhan/HarderNarasimhanProofs/FrameLinearAlgebraFirst.lean}{\texttt{frame\_linear\_algebra\_first}}.
Recibe los coeficientes positivos, $R>0$, las dos hipótesis de trazas
y la cota de codimensión total. Expresa la conclusión como la dimensión
del complemento ortogonal. La no negatividad de $\varepsilon$ se deduce
de las hipótesis; el caso $D=0$ se trata aparte. La prueba formal leída
sigue el cálculo de varianza de B. La prueba por $I-P$ desarrollada
arriba sigue A; las conclusiones y las constantes coinciden.

\item \textbf{La transferencia precisa.}
En el mismo archivo, \texttt{TheoremFSharp} define el enunciado
preciso, cuya prueba es
\href{https://github.com/CocoSgt/riemann-hypothesis-lean/blob/7cbf7a0c32fca1badf05931879a37c3f8a2b03b4/src/Formalism/HarderNarasimhan/HarderNarasimhanProofs/TheoremFSharp.lean}{\texttt{frame\_theorem\_sharp}}.
La interfaz pública
\href{https://github.com/CocoSgt/riemann-hypothesis-lean/blob/7cbf7a0c32fca1badf05931879a37c3f8a2b03b4/src/Formalism/HarderNarasimhan/ShiftInvariance.lean}{\texttt{theoremFSharp}}
la expone directamente. Sus hipótesis incluyen la integridad, las alturas
de los mapas no nulos, la media ponderada por $r_i$, los extremos de
pendiente, las dos cotas de marco y la relación de rangos.
El levantamiento se escribe por matrices: su bloque superior representa
el mapa aritmético después de cambios de coordenadas invertibles.
La proyección sobre ese bloque realiza la sobreyección de la prueba.
La interfaz admite fuentes individuales de rango cero si $R>0$;
nuestra exposición usa el caso de rangos positivos de A y B.

\item \textbf{La normalización y los subobjetos.}
\href{https://github.com/CocoSgt/riemann-hypothesis-lean/blob/7cbf7a0c32fca1badf05931879a37c3f8a2b03b4/src/Formalism/HarderNarasimhan/StatementPl.lean}{\texttt{StatementPl.lean}}
define \texttt{ht}, \texttt{size}, \texttt{muDef} y \texttt{htMap}.
La identidad demostrada por \texttt{size\_maximal\_slope} identifica
\texttt{muDef} con \eqref{eq:marcos-covolumen}. Internamente,
si $\lambda$ es una pendiente calculada mediante el covolumen euclídeo,
la pendiente de nuestra normalización es $2\lambda+c_K$, donde
\[
 c_K=-\log2+\frac{\log|\Delta_K|}{d_K}.
\]
Este término ya está incorporado en \texttt{muDef}.
Los extremos y los cortes recorren submódulos sobre $\OK$.
La realización sobre $\Z$ sirve para calcular rango y covolumen;
no amplía la familia de subobjetos a todos los subretículos sobre $\Z$.
Las multiplicidades de las listas formales usan
$\operatorname{rk}_{\Z}=d_K\operatorname{rk}_{\OK}$; el factor común
$d_K$ cancela en las medias y en las proporciones.

\item \textbf{Los cortes y su integral.}
En la prueba de transferencia, \texttt{thmF\_share} estima la
fracción de pendientes bajas y \texttt{thmF\_place} aplica el marco
a los cortes y compara los corangos tras proyectar.
El paso de altura a pendiente de la imagen utiliza
\texttt{slope\_inequality\_mean}; su prueba leída cambia métricas
para igualar las normas en los lugares y aplica una desigualdad
euclídea importada. La reconstrucción anterior usa determinantes
y saturación. Finalmente,
\href{https://github.com/CocoSgt/riemann-hypothesis-lean/blob/7cbf7a0c32fca1badf05931879a37c3f8a2b03b4/src/Formalism/CorankShare/CorankShareProofs/LayerCakeFormula.lean}{\texttt{layer\_cake\_formula}},
del módulo de corangos, integra los recuentos
estrictos, y \texttt{frameB\_integral} calcula el término logarítmico.
La prueba trata expresamente $u=0$ y los mapas nulos.
\end{enumerate}

Hay variantes que deben distinguirse: \texttt{FrameLA}, en el módulo
de corangos, tiene coeficiente $2\eta$; \texttt{TheoremF} resta además
$\log|I|$. Las declaraciones precisas para las fórmulas de esta nota
son \texttt{FrameLA1} y \texttt{TheoremFSharp}.

En este tramo no se encontró una discrepancia de signos o constantes
con A y B. Se distinguieron las definiciones de tipo \texttt{Prop}
de las pruebas leídas. El cotejo no desplegó toda la construcción de la
filtración ni todas las dependencias importadas, y no incluyó compilación,
ejecución de Lean o comparación integral con la distribución original.
Las hipótesis que debe producir la familia automorfa siguen siendo las
entradas de la subsección anterior.

\subsection{Ejercicios con soluciones}

\begin{ejercicio}\label{ej:marcos-proyeccion}
Sea $\pi:W\to V$ una aplicación lineal sobreyectiva entre espacios
vectoriales de dimensión finita, y sea $S\subset W$ un subespacio.
Demuestre que
$\operatorname{codim}_V\pi S\le\operatorname{codim}_W S$.
Determine cuándo hay igualdad.
\end{ejercicio}

\paragraph{Solución del Ejercicio~\ref{ej:marcos-proyeccion}.}
La fórmula de dimensiones da
\[
 \operatorname{codim}_V\pi S
 =\operatorname{codim}_W S-\dim\ker\pi+\dim(S\cap\ker\pi).
\]
El término restado es no negativo. Hay igualdad si y solo si
$\ker\pi\subseteq S$. El cálculo usa sobreyectividad y dimensiones,
sin recurrir a ninguna métrica.

\begin{ejercicio}\label{ej:marcos-pesos}
Dos fuentes tienen rangos $1$ y $9$, y pendientes desplazadas
constantes $9$ y $-1$, respectivamente. Calcule la media de la lista
completa y compárela con la media que asigna peso $1/2$ a cada fuente.
\end{ejercicio}

\paragraph{Solución del Ejercicio~\ref{ej:marcos-pesos}.}
La lista consta de un $9$ y nueve valores $-1$, así que su media es
$(9-9)/10=0$. La media de las dos fuentes con pesos iguales sería
$(9-1)/2=4$. El segundo número no es la media que controla la
codimensión total de los cortes: el marco cuenta las diez dimensiones.

\begin{ejercicio}\label{ej:marcos-altura-negativa}
Sobre $\Q$, sea $E=F=\Z$ como módulo, pero con normas
$\|1\|_E=1$ y $\|1\|_F=e^{-a/2}$, donde $a>0$.
Para la identidad integral $B:E\to F$, calcule $h(B)$ y las
pendientes en la normalización duplicada. Compruebe
\eqref{eq:marcos-morfismo-pendientes}.
\end{ejercicio}

\paragraph{Solución del Ejercicio~\ref{ej:marcos-altura-negativa}.}
La norma de operador es $e^{-a/2}$, por lo que $h(B)=-a$.
Se tiene $\mu(E)=0$ y $\mu(F)=a$; como los rangos son uno,
las pendientes máxima, mínima y media coinciden. La desigualdad
es $0\le a-a=0$. Este ejemplo conserva la integridad del mapa
y muestra por qué la altura arquimediana puede ser negativa.

\clearpage
\phantomsection
\addcontentsline{toc}{section}{Fuentes de esta entrega}
\renewcommand{\refname}{Fuentes de esta entrega}
\begin{thebibliography}{L}

\bibitem[A]{fuenteA}
\emph{Dirichlet L-functions and the critical line}.
Documento de 55 páginas, ejemplar suministrado para este estudio el
22 de noviembre de 2026 bajo el nombre
\texttt{riemann\_hypothesis.pdf}.
El ejemplar consultado no identifica un autor en la portada.
Las referencias internas A remiten a sus secciones, proposiciones,
ecuaciones y páginas; no atribuyen autoría al código ni al resumen.

\bibitem[B]{fuenteB}
Robin Zhang.
\emph{A preliminary note on Alpöge's proposed proof of the generalized
Riemann hypothesis}.
21 de noviembre de 2026, 73 páginas.
Ejemplar suministrado bajo el nombre \texttt{GRH note.pdf}.
La declaración del autor ocupa las pp.~1--2; la exposición matemática,
las pp.~3--67; las referencias comienzan en p.~67.

\bibitem[L]{lean}
Desarrollo Lean de Levent Alpöge. Código consultado mediante el
\href{https://github.com/CocoSgt/riemann-hypothesis-lean/tree/7cbf7a0c32fca1badf05931879a37c3f8a2b03b4}{espejo de CocoSgt}, revisión \texttt{7cbf7a0}.
El registro del proyecto conserva el identificador completo, las rutas
y los hashes de los archivos recuperados. La configuración declara
Lean~4.33.1 y la revisión de Mathlib identificada en ese registro.
El cotejo de esta entrega se limita a las interfaces y pasos señalados
en la Nota~\ref{sec:ventanas}. No incluyó compilación ni ejecución de
Lean, ni una comparación integral con la distribución original.

\end{thebibliography}

\subsection*{Estado de la entrega}

Se han desarrollado la transformada de Mellin de la suma suavizada,
las correcciones en $s=1$, el cálculo del residuo en un cero, el ejemplo
de proyección sobre $\Z^2$ y el lema local de denominadores. Se añadió
la demostración completa del criterio unilateral de ventanas: derivadas,
concentración, selección de polos, estimaciones de Cauchy y contradicción.
Se ha comprobado la correspondencia de estos pasos con A y B; para
el criterio de ventanas se leyeron además las declaraciones y pruebas
Lean identificadas en la nota y en el registro del proyecto.

Permanecen por estudiar la construcción de la combinación positiva
por caracteres, las
estimaciones de marcos y pendientes, y los resultados automorfos,
geométricos y cohomológicos que suministran el presupuesto final.
No se ha efectuado en esta entrega una auditoría del repositorio Lean.

\end{document}
